\documentclass[../main.tex]{subfiles}
\begin{document}

\chapter{分散共有メモリ}
\lessoninfo{
  video={P4L3},
  textbook={Nitzberg と Lo \cite{nitzberg1991}；Tanenbaum と Bos \cite{tanenbaum2015} ch.~8；Tanenbaum と van Steen \cite{tanenbaum2007} ch.~7},
  keywords={分散共有メモリ，RDMA，共有の粒度，偽共有，アクセスアルゴリズム，マイグレーション，レプリケーション，一貫性管理，ページフォールト，一貫性モデル，厳密一貫性，逐次一貫性，因果一貫性，弱一貫性},
}

第14章では分散ファイルシステムを扱った．そこではサーバがファイルを所有し，
その一部をキャッシュするクライアントにそれらを提供する．しかし多くの分散
アプリケーションでは，どのノードも状態の一部を所有し，かつそれを他の
ノードへ提供する．この章では，そのように分散した状態がどう管理されるかの
例として，分散共有メモリを調べる．分散共有メモリは，異なるマシン上の
プロセスが，あたかも単一の共有メモリ型マルチプロセッサ上で動作している
かのようにメモリを共有できるようにするものである．DSM システムの設計上の
選択---粒度，アクセスパターン，マイグレーションとレプリケーションの選択，
そして一貫性管理---と，DSM がページをどう探し出し MMU を使ってアクセスを
どう捕捉するか，そして，プロセスがメモリを共有したときに何を観測するかを
定める一貫性モデルを扱う．

\section{ピア分散アプリケーション}
\label{sec:l15-peer}

第14章の分散ファイルシステムでは，マシンの役割が明確に分かれている．
\source{P4L3，レッスンの概要，分散ファイルシステムの復習，ピア分散
アプリケーション；Tanenbaum と van Steen \cite{tanenbaum2007} ch.~2．}
サーバは状態，すなわちファイルを所有して管理し，サービス，すなわちファイル
アクセスを提供する．クライアントはサーバへ要求を送る．クライアントは状態の
一部をキャッシュし，それによって自身のアクセスのレイテンシが下がり，また
サーバはより多くのクライアントを支えられるようになるが，状態そのものは
サーバに属する．

多くの分散アプリケーションには，そのような分担が存在しない．ビッグデータ
解析，ウェブ検索，コンテンツ共有では，\margin{第16章では，こうした
アプリケーションをデータセンタで扱う．そこではフレームワークが，各計算を，
その入力をすでに保持しているノードへ送る．}%
アプリケーションの状態はその全ノードに
分散している．各ノードは状態の一部を所有し，その部分に対するサービスを他の
ノードへ提供すると同時に，他のノードが所有する状態を利用する．
したがってどのノードもクライアントでありサーバでもあり，すべてのノードが
ピアである．
このようなアプリケーションを\emph{ピア分散アプリケーション}と呼ぶ．ここで
「ピア」という語は緩い意味で使われている．ピアツーピアシステム（第14章）
\index{ぴあつーぴあ@ピアツーピア}では，システム全体の制御と管理までもが
すべてのピアによって行われるのに対し，ピア分散アプリケーションは，どの
ノードが状態のどの部分を保持しているかを追跡するといった管理作業を，
特定のノードに割り当ててよい．

ピアにまたがって広がる状態の管理は3つの問題を生じる．これらはこの章を
通じて繰り返し現れる．\margin{最初の2つは1台のマシンの中にも現れる．
NUMA 対応スケジューリング（第7章）は計算をそのメモリのもとへ動かす．}
\begin{description}
  \item[配置] 状態の各部分がどのノードに存在するか．最も頻繁にアクセス
    するプロセスのノードに状態が存在するとき，アクセスは最も速い．
  \item[マイグレーション] 他のノードに存在する状態を，移動によって，
    あるいは複製によって，それにアクセスするプロセスのノードへどのように，
    いつ持ってくるか．\caution{この章のマイグレーションはノード間で
    \emph{メモリ}を移動する．プロセス全体を別のマシンへ移す第8章のプロセス
    マイグレーションではない．}
  \item[共有の規則] 同じ状態に対する並行な操作がどの順序で効果を持つか，
    そして各ノードが他のノードの行った更新をいつ観測するかを定める規則．
\end{description}
状態がメモリであるピア分散アプリケーションである分散共有メモリが，この章で
これらの問題を調べるための例である．

\section{分散共有メモリ}
\label{sec:l15-dsm}

共有メモリ型マルチプロセッサ（第7章）%
\index{きょうゆうめもりがたまるちぷろせっさ@共有メモリ型マルチプロセッサ}%
は，異なる CPU 上のスレッドが，
そのすべてがアクセスするメモリを通じて協調することを可能にするが，その
メモリは1台のマシンが保持できる量に限られる．\source{P4L3，分散共有メモリ；
Nitzberg と Lo \cite{nitzberg1991}；Tanenbaum と Bos ch.~8．}
分散共有メモリは，このモデルを多数のマシンへ拡張する．

\begin{definition}[分散共有メモリ]
  ネットワークのノードにまたがってメモリを管理し，その上で動作する
  アプリケーションに単一の共有メモリマシン上で動作しているかのような錯覚を
  与えるサービスを
  \term[ぶんさんきょうゆうめもり]{分散共有メモリ}[distributed shared memory]%
  \index{DSM|see{分散共有メモリ}}（DSM）と呼ぶ．どのノードのプロセスも1つの
  共有アドレス空間を読み書きし，その内容は異なるノードの物理メモリに存在
  する．
\end{definition}

他のあらゆるピア分散アプリケーションと同様に，DSM システムの各ノードは
次のことを行う（\cref{fig:l15-dsm}）．
\begin{itemize}
  \item 状態，すなわちシステムの物理メモリの一部---自身のメモリのうち
    DSM に提供する部分---を所有する．
  \item その状態に対するサービス，すなわち自身のメモリに対する読み出しと
    書き込みを提供する．これらはどのノードのプロセスが発行してもよい．
  \item 他のノードとの一貫性プロトコルに参加する．これは，複数のノードが
    同じメモリにアクセスしたときにも共有メモリへのアクセスが意味を持ち
    続けること---たとえば読み出しが最も新しい書き込みの値を返すこと---を
    保証する．\caution{これは第9章の共有メモリ IPC ではない．そちらでは
    1台のマシンのメモリをプロセスどうしが共有し，その見え方はハードウェアが
    一貫させる．}
\end{itemize}

\begin{figure}[htbp]
  \centering
  % Distributed shared memory: processes on three nodes read and write one
% shared address space; every page of it resides in the memory that one of
% the nodes contributes, and remote accesses cross the network.
% Usage: % Distributed shared memory: processes on three nodes read and write one
% shared address space; every page of it resides in the memory that one of
% the nodes contributes, and remote accesses cross the network.
% Usage: % Distributed shared memory: processes on three nodes read and write one
% shared address space; every page of it resides in the memory that one of
% the nodes contributes, and remote accesses cross the network.
% Usage: \input{figures/l15-dsm}   (width ~118 mm)
\begin{tikzpicture}[x=1mm, y=1mm, font=\scriptsize\sffamily,
  >={Stealth[length=1.3mm]},
  mach/.style={draw, rounded corners=0.8mm, fill=white},
  proc/.style={draw, circle, fill=white, inner sep=0pt, minimum size=4.8mm},
  dsm/.style={draw, fill=ln-box, minimum width=30mm, minimum height=4.4mm, inner sep=0pt},
  frm/.style={draw, minimum width=8.4mm, minimum height=4.6mm, inner sep=0pt},
  cell/.style={draw, minimum width=13mm, minimum height=5mm, inner sep=0pt},
  lab/.style={font=\scriptsize\sffamily, text=ln-muted, inner sep=0.8pt},
  who/.style={font=\scriptsize\sffamily\bfseries, anchor=north west, inner sep=1pt}]
  \colorlet{lnNodeA}{ln-accent!22}
  \colorlet{lnNodeB}{ln-tip!22}
  \colorlet{lnNodeC}{ln-caution!16}
  % shared address space
  \node[lab, anchor=south west] at (1,40.8)
    {\iflangja{共有アドレス空間（すべてのプロセスから見える）}{shared address space (seen by every process)}};
  \foreach \ci/\cc in {0/lnNodeA,1/lnNodeB,2/lnNodeC,3/lnNodeA,4/lnNodeB,5/lnNodeC,6/lnNodeA,7/lnNodeB,8/lnNodeC} {
    \node[cell, fill=\cc] at ({7.5+13*\ci},38) {$p_{\ci}$};
  }
  % nodes
  \foreach \nk/\nx/\ncol/\npa/\npb/\fa/\fb/\fc in
      {1/0/lnNodeA/1/2/0/3/6, 2/41/lnNodeB/3/4/1/4/7, 3/82/lnNodeC/5/6/2/5/8} {
    \draw[mach] (\nx,0) rectangle ++(36,26);
    \node[who] at (\nx+0.5,25.5) {\iflangja{ノード}{node} \nk};
    \node[proc] at (\nx+15,19.5) {P\npa};
    \node[proc] at (\nx+27,19.5) {P\npb};
    \draw[<->] (\nx+15,21.9) -- (\nx+15,35.4);
    \draw[<->] (\nx+27,21.9) -- (\nx+27,35.4);
    \node[dsm] at (\nx+18,12) {DSM};
    \foreach \fj/\fp in {0/\fa,1/\fb,2/\fc} {
      \node[frm, fill=\ncol] at (\nx+8.5+9.5*\fj,4.6) {$p_{\fp}$};
    }
    \draw[thick, ln-muted] (\nx+18,0) -- (\nx+18,-5);
  }
  \node[lab, anchor=west] at (27.6,30.5) {\iflangja{読み書き}{read / write}};
  % network
  \draw[line width=1.2pt, ln-muted] (0,-5) -- (118,-5);
  \node[lab, anchor=north] at (59,-5.8)
    {\iflangja{ネットワーク（RDMA 対応のインターコネクトなど）}{network (e.g.\ an interconnect with RDMA)}};
\end{tikzpicture}
   (width ~118 mm)
\begin{tikzpicture}[x=1mm, y=1mm, font=\scriptsize\sffamily,
  >={Stealth[length=1.3mm]},
  mach/.style={draw, rounded corners=0.8mm, fill=white},
  proc/.style={draw, circle, fill=white, inner sep=0pt, minimum size=4.8mm},
  dsm/.style={draw, fill=ln-box, minimum width=30mm, minimum height=4.4mm, inner sep=0pt},
  frm/.style={draw, minimum width=8.4mm, minimum height=4.6mm, inner sep=0pt},
  cell/.style={draw, minimum width=13mm, minimum height=5mm, inner sep=0pt},
  lab/.style={font=\scriptsize\sffamily, text=ln-muted, inner sep=0.8pt},
  who/.style={font=\scriptsize\sffamily\bfseries, anchor=north west, inner sep=1pt}]
  \colorlet{lnNodeA}{ln-accent!22}
  \colorlet{lnNodeB}{ln-tip!22}
  \colorlet{lnNodeC}{ln-caution!16}
  % shared address space
  \node[lab, anchor=south west] at (1,40.8)
    {\iflangja{共有アドレス空間（すべてのプロセスから見える）}{shared address space (seen by every process)}};
  \foreach \ci/\cc in {0/lnNodeA,1/lnNodeB,2/lnNodeC,3/lnNodeA,4/lnNodeB,5/lnNodeC,6/lnNodeA,7/lnNodeB,8/lnNodeC} {
    \node[cell, fill=\cc] at ({7.5+13*\ci},38) {$p_{\ci}$};
  }
  % nodes
  \foreach \nk/\nx/\ncol/\npa/\npb/\fa/\fb/\fc in
      {1/0/lnNodeA/1/2/0/3/6, 2/41/lnNodeB/3/4/1/4/7, 3/82/lnNodeC/5/6/2/5/8} {
    \draw[mach] (\nx,0) rectangle ++(36,26);
    \node[who] at (\nx+0.5,25.5) {\iflangja{ノード}{node} \nk};
    \node[proc] at (\nx+15,19.5) {P\npa};
    \node[proc] at (\nx+27,19.5) {P\npb};
    \draw[<->] (\nx+15,21.9) -- (\nx+15,35.4);
    \draw[<->] (\nx+27,21.9) -- (\nx+27,35.4);
    \node[dsm] at (\nx+18,12) {DSM};
    \foreach \fj/\fp in {0/\fa,1/\fb,2/\fc} {
      \node[frm, fill=\ncol] at (\nx+8.5+9.5*\fj,4.6) {$p_{\fp}$};
    }
    \draw[thick, ln-muted] (\nx+18,0) -- (\nx+18,-5);
  }
  \node[lab, anchor=west] at (27.6,30.5) {\iflangja{読み書き}{read / write}};
  % network
  \draw[line width=1.2pt, ln-muted] (0,-5) -- (118,-5);
  \node[lab, anchor=north] at (59,-5.8)
    {\iflangja{ネットワーク（RDMA 対応のインターコネクトなど）}{network (e.g.\ an interconnect with RDMA)}};
\end{tikzpicture}
   (width ~118 mm)
\begin{tikzpicture}[x=1mm, y=1mm, font=\scriptsize\sffamily,
  >={Stealth[length=1.3mm]},
  mach/.style={draw, rounded corners=0.8mm, fill=white},
  proc/.style={draw, circle, fill=white, inner sep=0pt, minimum size=4.8mm},
  dsm/.style={draw, fill=ln-box, minimum width=30mm, minimum height=4.4mm, inner sep=0pt},
  frm/.style={draw, minimum width=8.4mm, minimum height=4.6mm, inner sep=0pt},
  cell/.style={draw, minimum width=13mm, minimum height=5mm, inner sep=0pt},
  lab/.style={font=\scriptsize\sffamily, text=ln-muted, inner sep=0.8pt},
  who/.style={font=\scriptsize\sffamily\bfseries, anchor=north west, inner sep=1pt}]
  \colorlet{lnNodeA}{ln-accent!22}
  \colorlet{lnNodeB}{ln-tip!22}
  \colorlet{lnNodeC}{ln-caution!16}
  % shared address space
  \node[lab, anchor=south west] at (1,40.8)
    {\iflangja{共有アドレス空間（すべてのプロセスから見える）}{shared address space (seen by every process)}};
  \foreach \ci/\cc in {0/lnNodeA,1/lnNodeB,2/lnNodeC,3/lnNodeA,4/lnNodeB,5/lnNodeC,6/lnNodeA,7/lnNodeB,8/lnNodeC} {
    \node[cell, fill=\cc] at ({7.5+13*\ci},38) {$p_{\ci}$};
  }
  % nodes
  \foreach \nk/\nx/\ncol/\npa/\npb/\fa/\fb/\fc in
      {1/0/lnNodeA/1/2/0/3/6, 2/41/lnNodeB/3/4/1/4/7, 3/82/lnNodeC/5/6/2/5/8} {
    \draw[mach] (\nx,0) rectangle ++(36,26);
    \node[who] at (\nx+0.5,25.5) {\iflangja{ノード}{node} \nk};
    \node[proc] at (\nx+15,19.5) {P\npa};
    \node[proc] at (\nx+27,19.5) {P\npb};
    \draw[<->] (\nx+15,21.9) -- (\nx+15,35.4);
    \draw[<->] (\nx+27,21.9) -- (\nx+27,35.4);
    \node[dsm] at (\nx+18,12) {DSM};
    \foreach \fj/\fp in {0/\fa,1/\fb,2/\fc} {
      \node[frm, fill=\ncol] at (\nx+8.5+9.5*\fj,4.6) {$p_{\fp}$};
    }
    \draw[thick, ln-muted] (\nx+18,0) -- (\nx+18,-5);
  }
  \node[lab, anchor=west] at (27.6,30.5) {\iflangja{読み書き}{read / write}};
  % network
  \draw[line width=1.2pt, ln-muted] (0,-5) -- (118,-5);
  \node[lab, anchor=north] at (59,-5.8)
    {\iflangja{ネットワーク（RDMA 対応のインターコネクトなど）}{network (e.g.\ an interconnect with RDMA)}};
\end{tikzpicture}

  \caption{分散共有メモリ．共有アドレス空間の各ページ $p_i$ は，いずれかの
    ノードが提供したメモリに存在する．各ノードの DSM 層は，ローカルな
    ページへのアクセスを処理し，リモートのページをネットワーク越しに取得
    する．}
  \label{fig:l15-dsm}
\end{figure}

DSM は，アプリケーションが1台のマシンのメモリの限界を超えて規模を拡大する
ことを可能にする．1台のマシンが保持できる以上のメモリを必要とする
アプリケーションは，それを他のマシンから得る．また DSM は共有メモリをより
低い費用で提供する．非常に大きなメモリを持つ1台のマシンは，中程度のメモリを
持つ複数のコモディティマシンよりも不釣り合いに高価だからである．その代価は
レイテンシである．\margin{ローカルメモリへのアクセスは
\SI{100}{\nano\second} 程度，データセンタのイーサネット越しの往復は数十
マイクロ秒であるのに対し，\SI{200}{\giga\bit\per\second} の InfiniBand
ポート越しの片側 RDMA 読み出しは約 \SI{1}{\micro\second} である（NVIDIA
ConnectX アダプタの仕様）．}他のノードに存在するメモリへのアクセスは
ネットワークを越えるので，ローカルメモリへのアクセスより桁違いに遅い．
コモディティのインターコネクトは，リモートダイレクトメモリアクセスを提供
することでこの代価を下げる．

\begin{definition}[リモートダイレクトメモリアクセス]
  \term{RDMA}[remote direct memory access]%
  \index{remote direct memory access|see{RDMA}}では，ネットワーク
  インタフェースが，自身のマシンのメモリとリモートマシンのメモリとの間で
  データを直接転送する．その際，リモートマシンの CPU もオペレーティング
  システムも転送に関与しない．
\end{definition}

RDMA は DMA\index{DMA}（第11章）をネットワーク越しに拡張したものである．
通常は OS バイパス\index{OSばいぱす@OSバイパス}（第11章）と組み合わされ，
どちらのマシンのカーネルもデータ経路に乗らないようにする．\margin{InfiniBand
は RDMA をネイティブに提供する．RoCE と iWARP はイーサネット上で運ばれる
RDMA プロトコルである．}

\section{ハードウェア DSM とソフトウェア DSM}
\label{sec:l15-hwsw}

DSM はハードウェアで提供することも，ソフトウェアで提供することも，両者の
組み合わせで提供することもできる．\source{P4L3，ハードウェア DSM 対
ソフトウェア DSM；Nitzberg と Lo \cite{nitzberg1991}．}
\begin{description}
  \item[ハードウェア支援 DSM] 専用のインターコネクトに依存する．各ノードの
    オペレーティングシステムには，自ノードのメモリよりはるかに大きい，複数の
    ノードにまたがる物理メモリを管理しているかのように見える．他の
    ノードに存在するアドレスへのアクセスは，ネットワークインタフェース
    カード（NIC）によってそのノードへのメッセージに変換され，NIC は，リモート
    ノードの CPU を関与させることなく，リモートメモリに対するアトミック操作を
    含めて共有メモリの管理のあらゆる側面に参加する．このようなハードウェアは
    高価であり，高性能計算のように性能が最優先されるハイエンドのシステムで
    用いられる．\margin{Dolphin のアダプタや Sequent の NUMA-Q で使われた
    Scalable Coherent Interface（IEEE 1596-1992）はそうしたインターコネクト
    であり，ノードのメモリをハードウェアで一貫させる．}
  \item[ソフトウェア支援 DSM] すべてをソフトウェアで行う．メモリを保持する
    ノードの特定，その取得，アクセスの検出，そして複製の一貫性の維持である．
    そのソフトウェアはオペレーティングシステムの一部か，プログラミング言語の
    ランタイムの一部である．前者はページを，後者はオブジェクトを共有する
    （\cref{sec:l15-granularity}）．
  \item[ハイブリッド DSM] 作業を両者で分担する．ハードウェア，とくに MMU
    は，あらゆるアクセスで起こらなければならないこと，すなわちアドレスの
    変換と，現在のマッピングが許可しないアクセスに対する例外の発生を担う．
    ソフトウェアは，システム全体の知識を要する判断---たとえばどのページを
    プリフェッチ\index{ぷりふぇっち@プリフェッチ}するか，すなわちアクセス
    される前に他のノードから取得するか---を担う．
\end{description}

\section{共有の粒度}
\label{sec:l15-granularity}

DSM の最初の設計上の決定は共有の単位である．すなわち，DSM がノード間で
転送し，複製し，全体として一貫させるメモリの単位である．\source{P4L3，
DSM の設計：共有の粒度；Nitzberg と Lo；Tanenbaum と Bos ch.~8．}共有
メモリ型マルチプロセッサでは，ハードウェアが CPU キャッシュを，数十バイトの
単位であるキャッシュラインの粒度で一貫させる（第10章）
\index{きゃっしゅいっかんせい@キャッシュ一貫性}．これが可能なのは，
ハードウェアがナノ秒のうちにキャッシュ間を調整するからである．DSM では
調整のたびにネットワーク越しのメッセージが必要になるので，このような
細かい粒度は費用がかかりすぎる．候補は次のものである．
\begin{description}
  \item[変数] 個々の変数を共有するのも細かすぎる．変数は数バイトしか
    占めないので，アクセスごとの調整に必要なメッセージが，データそのものの
    価値をはるかに超える代価となる．
  \item[ページ] ページ---普通は \SI{4}{\kibi\byte}---は，オペレーティング
    システムがすでに管理している単位である．\margin{\SI{2}{\mebi\byte} の
    ラージページ（第8章）なら一度に 512 個の基本ページを覆う．転送は少なく
    大きくなるが，偽共有もそれだけ増える．}ページに基づく DSM は，
    オペレーティングシステムが MMU と協調して実装する
    （\cref{sec:l15-implementing}）．アプリケーションからは透過であり，
    アプリケーションが書かれた言語にも依存しない．
  \item[オブジェクト] 言語ランタイムは，プログラムが生成するオブジェクトを
    知っている．ランタイムで実装されるオブジェクトに基づく DSM は，
    プログラムのデータ構造に対応する単位を共有し，オペレーティングシステムの
    変更を必要としないが，そのランタイムを使うプログラムだけがその恩恵を
    受ける．
\end{description}

粒度が粗いほど，転送や調整1回あたりの代価をより多くのデータで償却できるが，
同時に無関係なデータが同じ単位に置かれることにもなる．\caution{偽共有は
単一のマシンでも起こる．2つの変数が同じキャッシュラインを共有し，一貫性
プロトコル（第10章）がそのラインを往復させる場合である．\SI{4}{\kibi\byte}
のページは \SI{64}{\byte} のラインを 64 本含み，その移動1回はキャッシュ
ミスではなくネットワーク転送である．}

\begin{definition}[偽共有]
  共有されていないデータ---それぞれ別のノードのプロセスがアクセスする
  データ---が同じ共有の単位に存在し，その結果，DSM がそれらのデータが共有
  されているかのようにその単位へのアクセスを調整してしまう状況を
  \term[ぎきょうゆう]{偽共有}[false sharing]と呼ぶ．
\end{definition}

\begin{figure}[htbp]
  \centering
  % False sharing: (a) x and y on the same page p, written by processes on
% different nodes, make p move on every write; (b) on separate pages p and
% q, nothing is transferred.
% Usage: % False sharing: (a) x and y on the same page p, written by processes on
% different nodes, make p move on every write; (b) on separate pages p and
% q, nothing is transferred.
% Usage: % False sharing: (a) x and y on the same page p, written by processes on
% different nodes, make p move on every write; (b) on separate pages p and
% q, nothing is transferred.
% Usage: \input{figures/l15-false-sharing}   (width ~118 mm)
\begin{tikzpicture}[x=1mm, y=1mm, font=\scriptsize\sffamily,
  >={Stealth[length=1.4mm]},
  mach/.style={draw, rounded corners=0.8mm, fill=white},
  half/.style={draw, fill=ln-accent!20, minimum width=6mm, minimum height=5mm, inner sep=0pt, font=\scriptsize\ttfamily},
  empty/.style={draw, fill=white, minimum width=6mm, minimum height=5mm, inner sep=0pt},
  gone/.style={draw, densely dashed, ln-muted, minimum width=12mm, minimum height=5mm, inner sep=0pt},
  lab/.style={font=\scriptsize\sffamily, inner sep=0.8pt},
  who/.style={font=\scriptsize\sffamily\bfseries, anchor=north west, inner sep=1pt},
  ttl/.style={font=\footnotesize\sffamily\bfseries, anchor=west}]
  % ---------------- (a) same page
  \node[ttl] at (0,23) {\iflangja{(a) \texttt{x} と \texttt{y} が同じページ}{(a) \texttt{x} and \texttt{y} on one page}};
  \draw[mach] (0,0) rectangle (21,17);
  \node[who] at (0.5,16.5) {\iflangja{ノード 1}{node 1}};
  \node[lab] at (10.5,11) {\iflangja{\texttt{x} に書く}{writes \texttt{x}}};
  \node[half] at (7.5,4.5) {x};
  \node[half] at (13.5,4.5) {y};
  \node[lab, anchor=west] at (16.8,4.5) {$p$};
  \draw[mach] (35,0) rectangle (56,17);
  \node[who] at (35.5,16.5) {\iflangja{ノード 2}{node 2}};
  \node[lab] at (45.5,11) {\iflangja{\texttt{y} に書く}{writes \texttt{y}}};
  \node[gone] at (45.5,4.5) {};
  \draw[->, thick, ln-caution] (21,13) to[bend left=35] node[lab, above, text=ln-caution] {$p$} (35,13);
  \draw[->, thick, ln-caution] (35,5) to[bend left=35] node[lab, below, text=ln-caution] {$p$} (21,5);
  % ---------------- (b) separate pages
  \begin{scope}[shift={(62,0)}]
    \node[ttl] at (0,23) {\iflangja{(b) 別々のページ}{(b) separate pages}};
    \draw[mach] (0,0) rectangle (21,17);
    \node[who] at (0.5,16.5) {\iflangja{ノード 1}{node 1}};
    \node[lab] at (10.5,11) {\iflangja{\texttt{x} に書く}{writes \texttt{x}}};
    \node[half] at (7.5,4.5) {x};
    \node[empty] at (13.5,4.5) {};
    \node[lab, anchor=west] at (16.8,4.5) {$p$};
    \draw[mach] (35,0) rectangle (56,17);
    \node[who] at (35.5,16.5) {\iflangja{ノード 2}{node 2}};
    \node[lab] at (45.5,11) {\iflangja{\texttt{y} に書く}{writes \texttt{y}}};
    \node[half] at (42.5,4.5) {y};
    \node[empty] at (48.5,4.5) {};
    \node[lab, anchor=west] at (51.8,4.5) {$q$};
    \node[lab, text=ln-muted, align=center] at (28,9) {\iflangja{転送\\なし}{no\\transfers}};
  \end{scope}
\end{tikzpicture}
   (width ~118 mm)
\begin{tikzpicture}[x=1mm, y=1mm, font=\scriptsize\sffamily,
  >={Stealth[length=1.4mm]},
  mach/.style={draw, rounded corners=0.8mm, fill=white},
  half/.style={draw, fill=ln-accent!20, minimum width=6mm, minimum height=5mm, inner sep=0pt, font=\scriptsize\ttfamily},
  empty/.style={draw, fill=white, minimum width=6mm, minimum height=5mm, inner sep=0pt},
  gone/.style={draw, densely dashed, ln-muted, minimum width=12mm, minimum height=5mm, inner sep=0pt},
  lab/.style={font=\scriptsize\sffamily, inner sep=0.8pt},
  who/.style={font=\scriptsize\sffamily\bfseries, anchor=north west, inner sep=1pt},
  ttl/.style={font=\footnotesize\sffamily\bfseries, anchor=west}]
  % ---------------- (a) same page
  \node[ttl] at (0,23) {\iflangja{(a) \texttt{x} と \texttt{y} が同じページ}{(a) \texttt{x} and \texttt{y} on one page}};
  \draw[mach] (0,0) rectangle (21,17);
  \node[who] at (0.5,16.5) {\iflangja{ノード 1}{node 1}};
  \node[lab] at (10.5,11) {\iflangja{\texttt{x} に書く}{writes \texttt{x}}};
  \node[half] at (7.5,4.5) {x};
  \node[half] at (13.5,4.5) {y};
  \node[lab, anchor=west] at (16.8,4.5) {$p$};
  \draw[mach] (35,0) rectangle (56,17);
  \node[who] at (35.5,16.5) {\iflangja{ノード 2}{node 2}};
  \node[lab] at (45.5,11) {\iflangja{\texttt{y} に書く}{writes \texttt{y}}};
  \node[gone] at (45.5,4.5) {};
  \draw[->, thick, ln-caution] (21,13) to[bend left=35] node[lab, above, text=ln-caution] {$p$} (35,13);
  \draw[->, thick, ln-caution] (35,5) to[bend left=35] node[lab, below, text=ln-caution] {$p$} (21,5);
  % ---------------- (b) separate pages
  \begin{scope}[shift={(62,0)}]
    \node[ttl] at (0,23) {\iflangja{(b) 別々のページ}{(b) separate pages}};
    \draw[mach] (0,0) rectangle (21,17);
    \node[who] at (0.5,16.5) {\iflangja{ノード 1}{node 1}};
    \node[lab] at (10.5,11) {\iflangja{\texttt{x} に書く}{writes \texttt{x}}};
    \node[half] at (7.5,4.5) {x};
    \node[empty] at (13.5,4.5) {};
    \node[lab, anchor=west] at (16.8,4.5) {$p$};
    \draw[mach] (35,0) rectangle (56,17);
    \node[who] at (35.5,16.5) {\iflangja{ノード 2}{node 2}};
    \node[lab] at (45.5,11) {\iflangja{\texttt{y} に書く}{writes \texttt{y}}};
    \node[half] at (42.5,4.5) {y};
    \node[empty] at (48.5,4.5) {};
    \node[lab, anchor=west] at (51.8,4.5) {$q$};
    \node[lab, text=ln-muted, align=center] at (28,9) {\iflangja{転送\\なし}{no\\transfers}};
  \end{scope}
\end{tikzpicture}
   (width ~118 mm)
\begin{tikzpicture}[x=1mm, y=1mm, font=\scriptsize\sffamily,
  >={Stealth[length=1.4mm]},
  mach/.style={draw, rounded corners=0.8mm, fill=white},
  half/.style={draw, fill=ln-accent!20, minimum width=6mm, minimum height=5mm, inner sep=0pt, font=\scriptsize\ttfamily},
  empty/.style={draw, fill=white, minimum width=6mm, minimum height=5mm, inner sep=0pt},
  gone/.style={draw, densely dashed, ln-muted, minimum width=12mm, minimum height=5mm, inner sep=0pt},
  lab/.style={font=\scriptsize\sffamily, inner sep=0.8pt},
  who/.style={font=\scriptsize\sffamily\bfseries, anchor=north west, inner sep=1pt},
  ttl/.style={font=\footnotesize\sffamily\bfseries, anchor=west}]
  % ---------------- (a) same page
  \node[ttl] at (0,23) {\iflangja{(a) \texttt{x} と \texttt{y} が同じページ}{(a) \texttt{x} and \texttt{y} on one page}};
  \draw[mach] (0,0) rectangle (21,17);
  \node[who] at (0.5,16.5) {\iflangja{ノード 1}{node 1}};
  \node[lab] at (10.5,11) {\iflangja{\texttt{x} に書く}{writes \texttt{x}}};
  \node[half] at (7.5,4.5) {x};
  \node[half] at (13.5,4.5) {y};
  \node[lab, anchor=west] at (16.8,4.5) {$p$};
  \draw[mach] (35,0) rectangle (56,17);
  \node[who] at (35.5,16.5) {\iflangja{ノード 2}{node 2}};
  \node[lab] at (45.5,11) {\iflangja{\texttt{y} に書く}{writes \texttt{y}}};
  \node[gone] at (45.5,4.5) {};
  \draw[->, thick, ln-caution] (21,13) to[bend left=35] node[lab, above, text=ln-caution] {$p$} (35,13);
  \draw[->, thick, ln-caution] (35,5) to[bend left=35] node[lab, below, text=ln-caution] {$p$} (21,5);
  % ---------------- (b) separate pages
  \begin{scope}[shift={(62,0)}]
    \node[ttl] at (0,23) {\iflangja{(b) 別々のページ}{(b) separate pages}};
    \draw[mach] (0,0) rectangle (21,17);
    \node[who] at (0.5,16.5) {\iflangja{ノード 1}{node 1}};
    \node[lab] at (10.5,11) {\iflangja{\texttt{x} に書く}{writes \texttt{x}}};
    \node[half] at (7.5,4.5) {x};
    \node[empty] at (13.5,4.5) {};
    \node[lab, anchor=west] at (16.8,4.5) {$p$};
    \draw[mach] (35,0) rectangle (56,17);
    \node[who] at (35.5,16.5) {\iflangja{ノード 2}{node 2}};
    \node[lab] at (45.5,11) {\iflangja{\texttt{y} に書く}{writes \texttt{y}}};
    \node[half] at (42.5,4.5) {y};
    \node[empty] at (48.5,4.5) {};
    \node[lab, anchor=west] at (51.8,4.5) {$q$};
    \node[lab, text=ln-muted, align=center] at (28,9) {\iflangja{転送\\なし}{no\\transfers}};
  \end{scope}
\end{tikzpicture}

  \caption{偽共有．(a) ノード 1 が \texttt{x} に書き，ノード 2 が \texttt{y}
    に書く．両者は同じページ $p$ にあるので，ページが往復する．(b) 別々の
    ページ $p$ と $q$ にあれば，どちらのアクセスもローカルに留まる．}
  \label{fig:l15-false-sharing}
\end{figure}

\begin{example}
  変数 \texttt{x} と \texttt{y} が同じページ $p$ にあるとする．ノード~1 の
  プロセスは \texttt{x} だけに書き，ノード~2 のプロセスは \texttt{y} だけに
  書く．各プロセスは \SI{10}{\milli\second} ごとに自分の変数に書き，2つの
  プロセスの書き込みは交互に起こる．両プロセスはデータを何も共有していないが，
  ページに基づく DSM からは2つのノードがページ $p$ に書いているように見える
  （\cref{fig:l15-false-sharing}a）．どの書き込みの時点でも $p$ の最新の内容は
  他方のノードにあるので，$p$ は \SI{5}{\milli\second} ごと，毎秒 200 回
  転送され，$200\times\SI{4}{\kibi\byte} = \SI{800}{\kibi\byte}$ のネットワーク
  トラフィックが毎秒生じる．転送に \SI{0.5}{\milli\second} かかるとすれば，
  各プロセスは \SI{10}{\milli\second} ごとに \SI{0.5}{\milli\second}，自分の
  時間の \SI{5}{\percent} を待つことになる．\texttt{x} と \texttt{y} を
  別々のページに置けば（\cref{fig:l15-false-sharing}b），最初のアクセスの後は
  ページは転送されない．
\end{example}

偽共有は，異なるノードがアクセスするデータを異なる単位に配置することで
避けられる．アクセスパターンが分かっていれば，プログラマ，コンパイラ，
あるいはメモリアロケータがそれを行える．

\section{アクセスアルゴリズム}
\label{sec:l15-access}

DSM が何を提供しなければならないかは，その上で動作するアプリケーションの
アクセスアルゴリズムに依存する．\source{P4L3，DSM の設計：アクセス
アルゴリズム；Nitzberg と Lo．}アクセスアルゴリズムは，同じメモリの単位を
同時にいくつのノードが読み書きするかを表す．次の3つが区別される．
\begin{description}
  \item[SRSW] \term[たんいつりーだたんいつらいた]{単一リーダ単一ライタ}%
    [single reader/single writer]\index{SRSW|see{単一リーダ単一ライタ}}
    （SRSW）アクセスでは，ある時点で1つのノードだけが単位を読み書きする．
    DSM の主な仕事は，他のノードに存在するメモリをアクセス可能にすることで
    ある．2つのノードが同時に1つの単位を使うことはないので，調整すべき
    並行アクセスは存在しない．
  \item[MRSW] \term[ふくすうりーだたんいつらいた]{複数リーダ単一ライタ}%
    [multiple readers/single writer]\index{MRSW|see{複数リーダ単一ライタ}}
    （MRSW）アクセスでは，複数のノードが同時に単位を読んでよいが，書く
    ノードは1つだけである．DSM は，あらゆる読み出しが最も新しい書き込みの
    値を返すことを保証しなければならない．
  \item[MRMW] \term[ふくすうりーだふくすうらいた]{複数リーダ複数ライタ}%
    [multiple readers/multiple writers]%
    \index{MRMW|see{複数リーダ複数ライタ}}（MRMW）アクセスでは，複数の
    ノードが単位を読みも書きもしてよい．DSM はさらに，すべてのノードが同じ
    順序で書き込みを観測し，読み出しが最も新しい書き込みを返すように，
    書き込みを正しく順序づけなければならない．\caution{アクセスアルゴリズムは
    アプリケーションが何を行うかを述べるものであって，DSM が課す規律では
    ない．MRMW の下で DSM がすべきことは，並行なライタに対処することで
    あって，それを禁じることではない．}
\end{description}

\section{マイグレーションとレプリケーション}
\label{sec:l15-migration}

DSM の主要な性能指標はアクセスレイテンシ，すなわち共有メモリの読み出しや
書き込みにかかる時間である．\source{P4L3，DSM の設計：マイグレーション対
レプリケーション；Nitzberg と Lo．}リモートアクセスはローカルアクセスより
桁違いに遅いので，リモートアクセスの割合がわずかであっても平均は急激に
上がる．1万回に1回のリモートアクセスですでに平均は5割増しになり，100回に
1回では平均は約 50 倍になる．

\begin{example}
  \label{ex:l15-latency}
  ローカルメモリへのアクセスに \SI{100}{\nano\second} かかり，まず他の
  ノードから取得しなければならないページへのアクセスに
  \SI{0.5}{\milli\second} かかるとする．アクセスのうち割合 $f$ がリモート
  であるとき，平均アクセスレイテンシは次式で与えられる．
  \begin{equation}
    \bar{t} = (1-f)\times\SI{100}{\nano\second} + f\times\SI{0.5}{\milli\second} .
    \label{eq:l15-latency}
  \end{equation}
  $f=10^{-4}$ では $\bar{t} \approx \SI{100}{\nano\second} +
  \SI{50}{\nano\second} = \SI{150}{\nano\second}$ となり，ローカルメモリの
  レイテンシの 1.5 倍である．$f=10^{-2}$ では $\bar{t} =
  \SI{99}{\nano\second} + \SI{5000}{\nano\second} \approx
  \SI{5.1}{\micro\second}$ となり，約 51 倍である．
\end{example}

したがって低いレイテンシは局所性を要求する．プロセスがアクセスするメモリは，
そのプロセス自身のノードに存在すべきである．\margin{\cref{ex:l15-latency}
の \SI{0.5}{\milli\second} は控えめな値である．\SI{4}{\kibi\byte} のページを
回線へ送り出す時間は \SI{100}{\giga\bit\per\second} では
\SI{0.33}{\micro\second} にすぎない．フォールト，要求，プロトコルの
オーバーヘッドを加えても，そうしたインターコネクト越しの取得は数マイクロ秒
である．}メモリを，それにアクセスするプロセスのもとへ持ってくる技法は2つある
（\cref{fig:l15-migration}）．
\begin{description}
  \item[マイグレーション] DSM は各単位のコピーを1つだけ保ち，それをアクセス
    するノードへ移動する．コピーが1つしかないので，一貫させるべきものは何も
    ない．マイグレーションは SRSW アクセスに適する．そこでは，1つのノードが
    ある期間その単位を使い続けた後に，別のノードがそれを必要とする．
    マイグレーションは，異なるノードがアクセスするたびに単位を転送するので，
    ノードが短い間隔で交代する場合には代価が高く，また複数のノードに同時に
    サービスすることもできない．
  \item[レプリケーション] DSM は単位を複数のノードにコピーし，各ノードは
    自分のローカルな複製にアクセスする．キャッシングはレプリケーション
    \index{れぷりけーしょん@レプリケーション}の一形態である．レプリケー
    ションはより一般的な技法であり，MRSW アクセスや MRMW アクセスが要求する
    ように，複数のノードが同時に単位を読み，あるいは書くことさえ可能にする．
    その代価は一貫性管理（\cref{sec:l15-management}）である．1つの複製が
    書かれたとき，他の複製は無効化されるか更新されなければならない．この
    オーバーヘッドは複製の数とともに増え，DSM は各単位の複製の数を制限する
    ことでそれを抑えられる．\tip{メッセージを数えよ．複製が $n$ 個あれば，
    書き込み1回に $n-1$ 回の無効化がかかり，無効化された各ノードでの次の
    アクセスに取得がかかる．レプリケーションは，それによってローカルで済む
    読み出しの数がこれらを上回る間は引き合う．}
\end{description}

\begin{figure}[htbp]
  \centering
  % Migration and replication: (a) the single copy of page p moves to node B,
% which accesses it; (b) nodes B and C hold replicas of p, and a write by A
% makes the DSM invalidate them.
% Usage: % Migration and replication: (a) the single copy of page p moves to node B,
% which accesses it; (b) nodes B and C hold replicas of p, and a write by A
% makes the DSM invalidate them.
% Usage: % Migration and replication: (a) the single copy of page p moves to node B,
% which accesses it; (b) nodes B and C hold replicas of p, and a write by A
% makes the DSM invalidate them.
% Usage: \input{figures/l15-migration}   (width ~118 mm)
\begin{tikzpicture}[x=1mm, y=1mm, font=\scriptsize\sffamily,
  >={Stealth[length=1.4mm]},
  mach/.style={draw, rounded corners=0.8mm, fill=white},
  own/.style={draw, fill=ln-accent!40, minimum width=9mm, minimum height=5mm, inner sep=0pt},
  rep/.style={draw, fill=ln-accent!12, minimum width=9mm, minimum height=5mm, inner sep=0pt},
  gone/.style={draw, densely dashed, ln-muted, minimum width=9mm, minimum height=5mm, inner sep=0pt},
  lab/.style={font=\scriptsize\sffamily, inner sep=0.8pt},
  who/.style={font=\scriptsize\sffamily\bfseries, anchor=north west, inner sep=1pt},
  ttl/.style={font=\footnotesize\sffamily\bfseries, anchor=west}]
  % \lnnodes: three node boxes A, B, C at x = 0, 20, 40
  \def\lnnodes{%
    \foreach \nx/\nn in {0/A,20/B,40/C} {
      \draw[mach] (\nx,0) rectangle ++(16,13);
      \node[who] at (\nx+0.4,12.6) {\nn};
    }}
  % ---------------- (a) migration
  \node[ttl] at (0,26) {\iflangja{(a) マイグレーション}{(a) migration}};
  \lnnodes
  \node[gone] at (8,5) {};
  \node[own] at (28,5) {$p$};
  \draw[->, thick, ln-accent] (8,13) to[bend left=40] node[lab, above, text=ln-accent] {\iflangja{$p$ を移動}{move $p$}} (28,13);
  \node[lab, anchor=north] at (28,-1) {\iflangja{$p$ にアクセス}{accesses $p$}};
  % ---------------- (b) replication
  \begin{scope}[shift={(62,0)}]
    \node[ttl] at (0,26) {\iflangja{(b) レプリケーション}{(b) replication}};
    \lnnodes
    \node[own] at (8,5) {$p$};
    \node[rep] at (28,5) {$p$};
    \node[rep] at (48,5) {$p$};
    \draw[->, thick, ln-caution] (10,13) to[bend left=40] (28,13);
    \draw[->, thick, ln-caution] (6,13) to[bend left=38] node[lab, above=0.8mm, fill=white, text=ln-caution] {\iflangja{無効化}{invalidate}} (48,13);
    \node[lab, anchor=north] at (8,-1) {\iflangja{$p$ に書く}{writes $p$}};
    \node[lab, anchor=north] at (38,-1) {\iflangja{$p$ を読む}{read $p$}};
  \end{scope}
\end{tikzpicture}
   (width ~118 mm)
\begin{tikzpicture}[x=1mm, y=1mm, font=\scriptsize\sffamily,
  >={Stealth[length=1.4mm]},
  mach/.style={draw, rounded corners=0.8mm, fill=white},
  own/.style={draw, fill=ln-accent!40, minimum width=9mm, minimum height=5mm, inner sep=0pt},
  rep/.style={draw, fill=ln-accent!12, minimum width=9mm, minimum height=5mm, inner sep=0pt},
  gone/.style={draw, densely dashed, ln-muted, minimum width=9mm, minimum height=5mm, inner sep=0pt},
  lab/.style={font=\scriptsize\sffamily, inner sep=0.8pt},
  who/.style={font=\scriptsize\sffamily\bfseries, anchor=north west, inner sep=1pt},
  ttl/.style={font=\footnotesize\sffamily\bfseries, anchor=west}]
  % \lnnodes: three node boxes A, B, C at x = 0, 20, 40
  \def\lnnodes{%
    \foreach \nx/\nn in {0/A,20/B,40/C} {
      \draw[mach] (\nx,0) rectangle ++(16,13);
      \node[who] at (\nx+0.4,12.6) {\nn};
    }}
  % ---------------- (a) migration
  \node[ttl] at (0,26) {\iflangja{(a) マイグレーション}{(a) migration}};
  \lnnodes
  \node[gone] at (8,5) {};
  \node[own] at (28,5) {$p$};
  \draw[->, thick, ln-accent] (8,13) to[bend left=40] node[lab, above, text=ln-accent] {\iflangja{$p$ を移動}{move $p$}} (28,13);
  \node[lab, anchor=north] at (28,-1) {\iflangja{$p$ にアクセス}{accesses $p$}};
  % ---------------- (b) replication
  \begin{scope}[shift={(62,0)}]
    \node[ttl] at (0,26) {\iflangja{(b) レプリケーション}{(b) replication}};
    \lnnodes
    \node[own] at (8,5) {$p$};
    \node[rep] at (28,5) {$p$};
    \node[rep] at (48,5) {$p$};
    \draw[->, thick, ln-caution] (10,13) to[bend left=40] (28,13);
    \draw[->, thick, ln-caution] (6,13) to[bend left=38] node[lab, above=0.8mm, fill=white, text=ln-caution] {\iflangja{無効化}{invalidate}} (48,13);
    \node[lab, anchor=north] at (8,-1) {\iflangja{$p$ に書く}{writes $p$}};
    \node[lab, anchor=north] at (38,-1) {\iflangja{$p$ を読む}{read $p$}};
  \end{scope}
\end{tikzpicture}
   (width ~118 mm)
\begin{tikzpicture}[x=1mm, y=1mm, font=\scriptsize\sffamily,
  >={Stealth[length=1.4mm]},
  mach/.style={draw, rounded corners=0.8mm, fill=white},
  own/.style={draw, fill=ln-accent!40, minimum width=9mm, minimum height=5mm, inner sep=0pt},
  rep/.style={draw, fill=ln-accent!12, minimum width=9mm, minimum height=5mm, inner sep=0pt},
  gone/.style={draw, densely dashed, ln-muted, minimum width=9mm, minimum height=5mm, inner sep=0pt},
  lab/.style={font=\scriptsize\sffamily, inner sep=0.8pt},
  who/.style={font=\scriptsize\sffamily\bfseries, anchor=north west, inner sep=1pt},
  ttl/.style={font=\footnotesize\sffamily\bfseries, anchor=west}]
  % \lnnodes: three node boxes A, B, C at x = 0, 20, 40
  \def\lnnodes{%
    \foreach \nx/\nn in {0/A,20/B,40/C} {
      \draw[mach] (\nx,0) rectangle ++(16,13);
      \node[who] at (\nx+0.4,12.6) {\nn};
    }}
  % ---------------- (a) migration
  \node[ttl] at (0,26) {\iflangja{(a) マイグレーション}{(a) migration}};
  \lnnodes
  \node[gone] at (8,5) {};
  \node[own] at (28,5) {$p$};
  \draw[->, thick, ln-accent] (8,13) to[bend left=40] node[lab, above, text=ln-accent] {\iflangja{$p$ を移動}{move $p$}} (28,13);
  \node[lab, anchor=north] at (28,-1) {\iflangja{$p$ にアクセス}{accesses $p$}};
  % ---------------- (b) replication
  \begin{scope}[shift={(62,0)}]
    \node[ttl] at (0,26) {\iflangja{(b) レプリケーション}{(b) replication}};
    \lnnodes
    \node[own] at (8,5) {$p$};
    \node[rep] at (28,5) {$p$};
    \node[rep] at (48,5) {$p$};
    \draw[->, thick, ln-caution] (10,13) to[bend left=40] (28,13);
    \draw[->, thick, ln-caution] (6,13) to[bend left=38] node[lab, above=0.8mm, fill=white, text=ln-caution] {\iflangja{無効化}{invalidate}} (48,13);
    \node[lab, anchor=north] at (8,-1) {\iflangja{$p$ に書く}{writes $p$}};
    \node[lab, anchor=north] at (38,-1) {\iflangja{$p$ を読む}{read $p$}};
  \end{scope}
\end{tikzpicture}

  \caption{ページ $p$ を，それにアクセスするノードのもとへ持ってくる．
    (a) マイグレーションは唯一のコピーを移動する．(b) レプリケーションは
    複数のノードに複製を保ち，それらは $p$ をローカルに読む．所有者 A に
    よる書き込みは B と C の複製を無効化する（破線）．}
  \label{fig:l15-migration}
\end{figure}

\begin{example}
  ページ $p$ は初めノード~A に存在し，ノード A，B，C が
  \cref{tab:l15-migration} の順序でそれにアクセスする．ページの転送には
  \SI{0.5}{\milli\second}，無効化メッセージには \SI{0.05}{\milli\second}
  かかる．\margin{無効化はページの識別子---数十バイト---しか運ばないので，
  その代価はデータの代価ではなくメッセージのやり取りの代価である．}
  マイグレーションは，$p$ を保持していないノードがアクセスするたびに $p$ を
  移動する．7 回の転送，\SI{3.5}{\milli\second} である．レプリケーションでは，
  B と C は複製を一度取得し，ステップ~6 の A の書き込みがそれらを無効化する
  まではローカルに読む．4 回の転送と 2 回の無効化，\SI{2.1}{\milli\second}
  である．8 回のアクセスをすべて A が行っていたなら，どちらの技法も何も
  転送しなかっただろう．逆に A と B が交互に書いていたなら，書き込みのたびに
  他方のノードの複製が無効化され，次のアクセスのたびにそれを取得し直すことに
  なり，レプリケーションはマイグレーションと同じ回数の転送に加えて無効化の
  代価を払うことになる．
\end{example}

\begin{table}[htbp]
  \centering
  \caption{マイグレーションの下と，書き込み時に無効化するレプリケーションの
    下での，ページ $p$ へのアクセス．}
  \label{tab:l15-migration}
  \small
  \setlength{\tabcolsep}{4pt}
  \begin{tabular}{@{}rl ll ll@{}}
    \toprule
    & & \multicolumn{2}{c}{マイグレーション} &
      \multicolumn{2}{c@{}}{レプリケーション} \\
    \cmidrule(lr){3-4}\cmidrule(l){5-6}
    ステップ & アクセス & $p$ の位置 & 転送 & 複製の位置 & メッセージ \\
    \midrule
    0 & （初期状態） & A & --- & A & --- \\
    1 & A が書く & A & --- & A & --- \\
    2 & B が読む & B & A$\to$B & A, B & A から取得 \\
    3 & C が読む & C & B$\to$C & A, B, C & A から取得 \\
    4 & B が読む & B & C$\to$B & A, B, C & --- \\
    5 & C が読む & C & B$\to$C & A, B, C & --- \\
    6 & A が書く & A & C$\to$A & A & B, C を無効化 \\
    7 & B が読む & B & A$\to$B & A, B & A から取得 \\
    8 & C が読む & C & B$\to$C & A, B, C & A から取得 \\
    \midrule
    \multicolumn{2}{@{}l}{合計} & & 7 回の転送 & &
      4 回の取得，2 回の無効化 \\
    \bottomrule
  \end{tabular}
\end{table}

どちらの技法もあらゆるワークロードに最良というわけではない．マイグレー
ションは，1つの単位が一度に1つのノードによって長い期間使われるときに適する．
レプリケーションは，複数のノードが同時に単位にアクセスするときに必要であり，
読み出しが支配的なときに引き合う．\margin{書き込み時に無効化する
レプリケーションは，第10章の書き込み無効化戦略をソフトウェアで行うもので
ある．キャッシュラインの代わりにページを，バストランザクションの代わりに
メッセージを用いる．}複製された単位が頻繁に書かれるときは，
書き込みのたびに起こる一貫性管理が，複製の局所性によって節約される分より
高くつきうる．

\begin{keypoint}
  DSM の性能はアクセスレイテンシによって決まり，低いレイテンシは局所性を
  要求する．マイグレーションは唯一のコピーを移動し，SRSW アクセスに適する．
  レプリケーションは複数のコピーを保ち，同時アクセスを支えるが，その代価と
  して，複製の数と書き込みの頻度とともに増える一貫性管理を払う．
\end{keypoint}

\section{一貫性管理}
\label{sec:l15-management}

いったんメモリが複製されると，DSM は複製を一貫させなければならない．
\source{P4L3，DSM の設計：一貫性管理；Nitzberg と Lo．}共有メモリ型
マルチプロセッサでは，ソフトウェアがロックで自身のアクセスを順序づけ，
ハードウェアが CPU キャッシュ内のメモリの複製を一貫させる（第10章）．DSM には
そのようなハードウェアがない．その機構は，分散ファイルシステムがクライアントの
キャッシュを一貫させる機構（第14章）に似ている．
\begin{description}
  \item[プッシュ] 複製が書かれたとき，DSM は同じ単位の他の複製を保持する
    ノードへ，直ちに無効化を押し出す．この方式は\emph{先回り的}（proactive）
    かつ\emph{即時的}（eager）である．また\emph{悲観的}（pessimistic）でも
    ある．古い複製は読まれて害をなすと仮定し，古い複製が生じ次第それを取り
    除くからである．\caution{DSM における一貫性管理はキャッシュ一貫性では
    ない．ソフトウェアによって，ページやオブジェクトの粒度で，しかも一貫性
    モデルが要求する時点でのみ行われる．}
  \item[プル] ノードは，単位の所有者から変更に関する情報を，周期的に，
    あるいは必要になったときに引き出す．この方式は\emph{反応的}（reactive，
    オンデマンド）かつ\emph{遅延的}（lazy）である．また\emph{楽観的}
    （optimistic）でもある．古い複製をしばらく使っても害はないと期待し，
    即時的な伝播に要するメッセージを節約するからである．
\end{description}
DSM が実際にこれらの機構を起動するのがいつか---書き込みのたびか，アクセスの
たびか，周期的にか，あるいはアプリケーションが指定する時点でのみか---は，
DSM が共有状態に対して提供する一貫性モデルによって決まる
（\cref{sec:l15-models}）．\margin{プッシュとプルは，第14章のサーバ駆動
およびクライアント駆動のキャッシュ一貫性に対応する DSM 版である．}

\begin{figure}[htbp]
  \centering
  % Consistency management: node A writes page p at 0.5 s, node B holds a
% replica and reads p at 1 s and 3 s.  (a) Push: A invalidates B's replica
% at once.  (b) Pull: B polls A every 2 s, and its replica is stale in
% between.
% Usage: % Consistency management: node A writes page p at 0.5 s, node B holds a
% replica and reads p at 1 s and 3 s.  (a) Push: A invalidates B's replica
% at once.  (b) Pull: B polls A every 2 s, and its replica is stale in
% between.
% Usage: % Consistency management: node A writes page p at 0.5 s, node B holds a
% replica and reads p at 1 s and 3 s.  (a) Push: A invalidates B's replica
% at once.  (b) Pull: B polls A every 2 s, and its replica is stale in
% between.
% Usage: \input{figures/l15-push-pull}   (width ~116 mm)
\begin{tikzpicture}[x=1mm, y=1mm, font=\scriptsize\sffamily,
  >={Stealth[length=1.3mm]},
  row/.style={font=\scriptsize\sffamily, anchor=east},
  lab/.style={font=\scriptsize\sffamily, inner sep=0.8pt},
  ev/.style={circle, fill=black, inner sep=0pt, minimum size=1.3mm},
  ttl/.style={font=\footnotesize\sffamily\bfseries, anchor=west}]
  % time t (s) -> x = 22 + 23 t (mm); rows: A at 9, B at 0 (inside a scope)
  \def\lnrows{%
    \node[row] at (20,9) {\iflangja{A（書き手）}{A (writer)}};
    \node[row] at (20,0) {\iflangja{B（複製）}{B (replica)}};
    \foreach \y in {0,9} { \draw[ln-rule] (22,\y) -- (116,\y); }
    \node[ev] at ({22+23*0.5},9) {};
    \node[lab, anchor=south] at ({22+23*0.5},10) {\iflangja{書く}{write}};
    \foreach \t in {1,3} { \node[ev] at ({22+23*\t},0) {}; }
  }
  % ---------------- (a) push
  \begin{scope}[shift={(0,27)}]
    \node[ttl] at (0,16) {\iflangja{(a) プッシュ}{(a) push}};
    \lnrows
    \draw[->, thick, ln-caution] ({22+23*0.5},8) -- ({22+23*0.5},1);
    \node[lab, anchor=east, text=ln-caution] at ({21.4+23*0.5},4.5) {\iflangja{無効化}{invalidate}};
    \draw[->, densely dashed] ({22+23*1},8) -- ({22+23*1},1);
    \node[lab, anchor=west] at ({22.6+23*1},4.5) {\iflangja{取得}{fetch}};
    \node[lab, anchor=north] at ({22+23*1},-1) {\iflangja{読む：新}{read: new}};
    \node[lab, anchor=north] at ({22+23*3},-1) {\iflangja{読む：新}{read: new}};
  \end{scope}
  % ---------------- (b) pull
  \begin{scope}[shift={(0,0)}]
    \node[ttl] at (0,16) {\iflangja{(b) プル（2 秒ごと）}{(b) pull (every 2 s)}};
    \lnrows
    \fill[ln-caution!22] ({22+23*0.5},-1) rectangle ({22+23*2},1);
    \draw[ln-rule] ({22+23*0.5},0) -- ({22+23*2},0);
    \foreach \t in {1,3} { \node[ev] at ({22+23*\t},0) {}; }
    \node[lab, anchor=south, text=ln-caution] at ({22+23*1.5},1.2) {\iflangja{古い}{stale}};
    \foreach \t in {0,2,4} {
      \draw[<->, densely dashed] ({22+23*\t},1) -- ({22+23*\t},8);
    }
    \node[lab, anchor=west] at ({22.6+23*2},4.5) {\iflangja{ポーリング}{poll}};
    \node[lab, anchor=north] at ({22+23*1},-1) {\iflangja{読む：古}{read: stale}};
    \node[lab, anchor=north] at ({22+23*3},-1) {\iflangja{読む：新}{read: new}};
  \end{scope}
  % time axis
  \draw[->] (22,-8) -- (117,-8);
  \foreach \t in {0,1,2,3,4} {
    \draw ({22+23*\t},-7.2) -- ({22+23*\t},-8.8);
    \node[lab, anchor=north] at ({22+23*\t},-9) {\t};
  }
  \node[lab, anchor=north east] at (20,-8.6) {\iflangja{時間（秒）}{time (s)}};
\end{tikzpicture}
   (width ~116 mm)
\begin{tikzpicture}[x=1mm, y=1mm, font=\scriptsize\sffamily,
  >={Stealth[length=1.3mm]},
  row/.style={font=\scriptsize\sffamily, anchor=east},
  lab/.style={font=\scriptsize\sffamily, inner sep=0.8pt},
  ev/.style={circle, fill=black, inner sep=0pt, minimum size=1.3mm},
  ttl/.style={font=\footnotesize\sffamily\bfseries, anchor=west}]
  % time t (s) -> x = 22 + 23 t (mm); rows: A at 9, B at 0 (inside a scope)
  \def\lnrows{%
    \node[row] at (20,9) {\iflangja{A（書き手）}{A (writer)}};
    \node[row] at (20,0) {\iflangja{B（複製）}{B (replica)}};
    \foreach \y in {0,9} { \draw[ln-rule] (22,\y) -- (116,\y); }
    \node[ev] at ({22+23*0.5},9) {};
    \node[lab, anchor=south] at ({22+23*0.5},10) {\iflangja{書く}{write}};
    \foreach \t in {1,3} { \node[ev] at ({22+23*\t},0) {}; }
  }
  % ---------------- (a) push
  \begin{scope}[shift={(0,27)}]
    \node[ttl] at (0,16) {\iflangja{(a) プッシュ}{(a) push}};
    \lnrows
    \draw[->, thick, ln-caution] ({22+23*0.5},8) -- ({22+23*0.5},1);
    \node[lab, anchor=east, text=ln-caution] at ({21.4+23*0.5},4.5) {\iflangja{無効化}{invalidate}};
    \draw[->, densely dashed] ({22+23*1},8) -- ({22+23*1},1);
    \node[lab, anchor=west] at ({22.6+23*1},4.5) {\iflangja{取得}{fetch}};
    \node[lab, anchor=north] at ({22+23*1},-1) {\iflangja{読む：新}{read: new}};
    \node[lab, anchor=north] at ({22+23*3},-1) {\iflangja{読む：新}{read: new}};
  \end{scope}
  % ---------------- (b) pull
  \begin{scope}[shift={(0,0)}]
    \node[ttl] at (0,16) {\iflangja{(b) プル（2 秒ごと）}{(b) pull (every 2 s)}};
    \lnrows
    \fill[ln-caution!22] ({22+23*0.5},-1) rectangle ({22+23*2},1);
    \draw[ln-rule] ({22+23*0.5},0) -- ({22+23*2},0);
    \foreach \t in {1,3} { \node[ev] at ({22+23*\t},0) {}; }
    \node[lab, anchor=south, text=ln-caution] at ({22+23*1.5},1.2) {\iflangja{古い}{stale}};
    \foreach \t in {0,2,4} {
      \draw[<->, densely dashed] ({22+23*\t},1) -- ({22+23*\t},8);
    }
    \node[lab, anchor=west] at ({22.6+23*2},4.5) {\iflangja{ポーリング}{poll}};
    \node[lab, anchor=north] at ({22+23*1},-1) {\iflangja{読む：古}{read: stale}};
    \node[lab, anchor=north] at ({22+23*3},-1) {\iflangja{読む：新}{read: new}};
  \end{scope}
  % time axis
  \draw[->] (22,-8) -- (117,-8);
  \foreach \t in {0,1,2,3,4} {
    \draw ({22+23*\t},-7.2) -- ({22+23*\t},-8.8);
    \node[lab, anchor=north] at ({22+23*\t},-9) {\t};
  }
  \node[lab, anchor=north east] at (20,-8.6) {\iflangja{時間（秒）}{time (s)}};
\end{tikzpicture}
   (width ~116 mm)
\begin{tikzpicture}[x=1mm, y=1mm, font=\scriptsize\sffamily,
  >={Stealth[length=1.3mm]},
  row/.style={font=\scriptsize\sffamily, anchor=east},
  lab/.style={font=\scriptsize\sffamily, inner sep=0.8pt},
  ev/.style={circle, fill=black, inner sep=0pt, minimum size=1.3mm},
  ttl/.style={font=\footnotesize\sffamily\bfseries, anchor=west}]
  % time t (s) -> x = 22 + 23 t (mm); rows: A at 9, B at 0 (inside a scope)
  \def\lnrows{%
    \node[row] at (20,9) {\iflangja{A（書き手）}{A (writer)}};
    \node[row] at (20,0) {\iflangja{B（複製）}{B (replica)}};
    \foreach \y in {0,9} { \draw[ln-rule] (22,\y) -- (116,\y); }
    \node[ev] at ({22+23*0.5},9) {};
    \node[lab, anchor=south] at ({22+23*0.5},10) {\iflangja{書く}{write}};
    \foreach \t in {1,3} { \node[ev] at ({22+23*\t},0) {}; }
  }
  % ---------------- (a) push
  \begin{scope}[shift={(0,27)}]
    \node[ttl] at (0,16) {\iflangja{(a) プッシュ}{(a) push}};
    \lnrows
    \draw[->, thick, ln-caution] ({22+23*0.5},8) -- ({22+23*0.5},1);
    \node[lab, anchor=east, text=ln-caution] at ({21.4+23*0.5},4.5) {\iflangja{無効化}{invalidate}};
    \draw[->, densely dashed] ({22+23*1},8) -- ({22+23*1},1);
    \node[lab, anchor=west] at ({22.6+23*1},4.5) {\iflangja{取得}{fetch}};
    \node[lab, anchor=north] at ({22+23*1},-1) {\iflangja{読む：新}{read: new}};
    \node[lab, anchor=north] at ({22+23*3},-1) {\iflangja{読む：新}{read: new}};
  \end{scope}
  % ---------------- (b) pull
  \begin{scope}[shift={(0,0)}]
    \node[ttl] at (0,16) {\iflangja{(b) プル（2 秒ごと）}{(b) pull (every 2 s)}};
    \lnrows
    \fill[ln-caution!22] ({22+23*0.5},-1) rectangle ({22+23*2},1);
    \draw[ln-rule] ({22+23*0.5},0) -- ({22+23*2},0);
    \foreach \t in {1,3} { \node[ev] at ({22+23*\t},0) {}; }
    \node[lab, anchor=south, text=ln-caution] at ({22+23*1.5},1.2) {\iflangja{古い}{stale}};
    \foreach \t in {0,2,4} {
      \draw[<->, densely dashed] ({22+23*\t},1) -- ({22+23*\t},8);
    }
    \node[lab, anchor=west] at ({22.6+23*2},4.5) {\iflangja{ポーリング}{poll}};
    \node[lab, anchor=north] at ({22+23*1},-1) {\iflangja{読む：古}{read: stale}};
    \node[lab, anchor=north] at ({22+23*3},-1) {\iflangja{読む：新}{read: new}};
  \end{scope}
  % time axis
  \draw[->] (22,-8) -- (117,-8);
  \foreach \t in {0,1,2,3,4} {
    \draw ({22+23*\t},-7.2) -- ({22+23*\t},-8.8);
    \node[lab, anchor=north] at ({22+23*\t},-9) {\t};
  }
  \node[lab, anchor=north east] at (20,-8.6) {\iflangja{時間（秒）}{time (s)}};
\end{tikzpicture}

  \caption{ノード A は \SI{0.5}{\second} にページ $p$ に書き，ノード B は
    複製を保持している．(a) A が無効化を押し出し，B は次の読み出しで新しい
    内容を取得する．(b) B は \SI{2}{\second} ごとにポーリングし，その複製は
    次のポーリングまで古いままである．}
  \label{fig:l15-push-pull}
\end{figure}

\begin{example}
  ノード~A がページ $p$ を所有し，ノード~B が複製を保持している．A は
  $t=\SI{0.5}{\second}$ に $p$ に書き，B は \SI{1}{\second} と
  \SI{3}{\second} に $p$ を読む（\cref{fig:l15-push-pull}）．プッシュでは，
  A は \SI{0.5}{\second} に B へ無効化を送る．B の \SI{1}{\second} の
  読み出しは有効な複製を見出せず，新しい内容を取得する．プルでは，B は
  \SI{2}{\second} ごとに $p$ が変更されたかを A に尋ねる．B の
  \SI{1}{\second} の読み出しは古い内容を返し，\SI{2}{\second} のポーリングが
  変更を明らかにし，\SI{3}{\second} の読み出しが新しい内容を返す．このように
  複製は最大で1ポーリング周期のあいだ古いままでありうる．その代わり，プルは
  一定の頻度でメッセージを要する．\SI{2}{\second} ごとにポーリングする
  10 個の複製は毎秒 5 個のメッセージを送る．$p$ が毎秒 100 回書かれようと，
  まったく書かれまいと同じである．プッシュは，$p$ が書かれない間はメッセージを
  送らず，書き込み1回につき複製ごとに高々1回の無効化を送る．\tip{プッシュは
  書き込みに比例して，プルは時間に比例してメッセージを要する．どちらが安いか
  は，書き込みの頻度とポーリングの頻度を比べればよい．}
\end{example}

\section{DSM のアーキテクチャ}
\label{sec:l15-architecture}

ここでは，オペレーティングシステムが実装するページに基づく DSM を考える．
\source{P4L3，DSM のアーキテクチャ，分散した状態の索引付け；Nitzberg と
Lo．}その構成は，これまでの設計上の選択から導かれる．
\begin{itemize}
  \item どのノードも主記憶の一部，いくつかのページフレームを DSM に提供
    する．提供されたページはシステム全体で1つの共有メモリのプールをなし，
    すべてのノードからアクセスできる．
  \item どのノードも，自分のプロセスがアクセスするページのローカルなコピー，
    すなわち複製のキャッシュを保持し，アクセスレイテンシを下げる．
  \item 管理もまた分散される．どのノードも共有メモリの一部について責任を
    持つ．
\end{itemize}
各ページについて，DSM は次の役割を区別する．
\begin{description}
  \item[ホームノード] ページを管理するノードであり，その\emph{管理ノード}
    とも呼ばれる．ホームノードはページの状態，すなわちどのノードがそれに
    アクセスしているか，変更されているか，そのキャッシングが有効か，
    ロックされているか，そして複製がどこにあるかを保持する．ページのアドレスは
    それを初めに管理するノードを名指す．管理はアドレスを変えずに他のノードへ
    引き渡せる（\cref{sec:l15-indexing}）．\margin{ホームノードは，分散
    ファイルシステム（第14章）でサーバが果たす役割を果たす．ただし，どの
    ノードもメモリの一部についてのホームノードである点が異なる．}
  \item[所有者] 現在，ページの正本のコピーを保持し，その更新を制御する
    ノード．たとえば，直近にそれを書く権利を得たノードである．所有権は
    ページがアクセスされるにつれて変わるので，所有者はホームノードである
    必要はない．ホームノードが現在の所有者を追跡する．
  \item[明示的な複製] アクセスの副作用としてではなく，DSM が意図的に作る
    複製．負荷分散のため，性能のため，あるいは信頼性のためであり，たとえば
    データセンタではデータを3つのコピーで保持するのが普通である．
    \margin{3つというのは HDFS の既定のレプリケーション数であり，HDFS は
    2つの複製を同じラックに，3つ目を別のラックに置く（Hadoop の
    ドキュメント）．}%
    これらの複製をどこに置き，どう管理するかはホームノードが制御する．
\end{description}

\subsection{分散した状態の索引付け}
\label{sec:l15-indexing}

アクセスを処理するために，ノードはページのメタデータを見つけなければ
ならない．メタデータはそのページの管理ノードが保持している．DSM はこの
メタデータを次のように組織する（\cref{fig:l15-indexing}）．
\begin{itemize}
  \item どのページも，ノード ID と物理フレーム番号
    \index{ぶつりふれーむばんごう@物理フレーム番号}（PFN，第8章）からなる
    アドレスで識別される．ノード ID は，そのフレームを提供し，初めはその
    ページを管理するノードを指す．PFN はそのノード上のフレームを指す．
  \item ページ ID から管理ノードの ID へのグローバルマップがすべての
    ノードに複製されるので，どのノードも任意のページの管理ノードを，
    メッセージを交換せずにローカルに見つけられる．
  \item ページ自体のメタデータは分割される．各管理ノードは自分が管理する
    ページのメタデータを保持する．\margin{マップを複製できるのは，それが
    小さいからである．後述の例の 8 ビットのノード ID であればエントリは
    256 個しかない．はるかに数の多いページごとのメタデータは分割されたまま
    である．}
\end{itemize}
ページ ID のビットから管理ノードを直接導く代わりに，グローバルマップを
\emph{マッピングテーブル}として実装することもできる．ページ ID の一部が
テーブルへの索引となり，そのエントリが管理ノードを示す．こうすれば，負荷
分散のためや故障の後に，すべてのノードでテーブルのエントリを更新するだけで，
あるページ集合の管理ノードを変更できる．ページ ID も，それを使うあらゆる参照も
変わらない．いったんテーブルを介在させれば，ノード ID のフィールドはテーブル
への索引にすぎない．ページの管理ノードは
エントリが示すノードであり，もはやそのフィールドが指すノードである必要は
ない．\tip{マッピングテーブルとページの管理ノードの関係は，ページテーブルと
フレームの関係と同じである．すなわち，名前を変えないまま名前の指す先を
変えられるようにする間接の層である．}

\begin{figure}[htbp]
  \centering
  % Indexing distributed state: the node-ID field of a page ID indexes the
% global mapping table, replicated on every node, which names the manager;
% the manager keeps the metadata of its pages.
% Usage: % Indexing distributed state: the node-ID field of a page ID indexes the
% global mapping table, replicated on every node, which names the manager;
% the manager keeps the metadata of its pages.
% Usage: % Indexing distributed state: the node-ID field of a page ID indexes the
% global mapping table, replicated on every node, which names the manager;
% the manager keeps the metadata of its pages.
% Usage: \input{figures/l15-indexing}   (width ~118 mm)
\begin{tikzpicture}[x=1mm, y=1mm, font=\scriptsize\sffamily,
  >={Stealth[length=1.4mm]},
  fld/.style={draw, minimum height=6mm, inner sep=0pt, font=\scriptsize\ttfamily},
  cellc/.style={draw, minimum height=5mm, inner sep=0pt},
  mach/.style={draw, rounded corners=0.8mm, fill=white},
  lab/.style={font=\scriptsize\sffamily, inner sep=0.8pt},
  hd/.style={font=\scriptsize\sffamily\bfseries, inner sep=0.8pt}]
  % page ID
  \node[hd, anchor=south west] at (0,37) {\iflangja{ページ ID}{page ID}};
  \node[fld, fill=ln-accent!20, minimum width=12mm] (nid) at (6,33) {05};
  \node[fld, fill=white, minimum width=22mm] at (23,33) {0012A3};
  \node[lab, anchor=north, text=ln-muted] at (6,29.6) {\iflangja{ノード ID}{node ID}};
  \node[lab, anchor=north, text=ln-muted] at (23,29.6) {PFN};
  % global mapping table
  \node[hd, anchor=south] at (58,37) {\iflangja{グローバルマップ}{global mapping table}};
  \node[lab, anchor=south, text=ln-muted] at (58,34) {\iflangja{（全ノードに複製）}{(replicated on every node)}};
  \node[lab, text=ln-muted] at (51,30.5) {\iflangja{索引}{index}};
  \node[lab, text=ln-muted] at (63,30.5) {\iflangja{管理ノード}{manager}};
  \foreach \ry/\ri/\rm in {26/04/4, 21/05/9, 16/06/6} {
    \node[cellc, minimum width=10mm, fill=white] at (51,\ry) {\ttfamily\ri};
    \node[cellc, minimum width=14mm, fill=white] at (63,\ry) {\rm};
  }
  \node[cellc, minimum width=14mm, fill=ln-accent!20] at (63,21) {9};
  \node[lab, text=ln-muted] at (51,11) {$\vdots$};
  \node[lab, text=ln-muted] at (63,11) {$\vdots$};
  \draw[->] (nid.south) ++(0,-3.6) -- (6,21) -- (46,21);
  % manager node
  \draw[mach] (78,4) rectangle (118,37);
  \node[hd, anchor=north west] at (78.5,36.5) {\iflangja{ノード 9（管理ノード）}{node 9 (manager)}};
  \node[lab, anchor=north west, text=ln-muted] at (78.5,32) {\iflangja{管理するページのメタデータ}{metadata of its pages}};
  \foreach \cx/\cw/\ct in {86.5/13/{\iflangja{ページ}{page}}, 99/12/{\iflangja{所有者}{owner}}, 110.5/11/{\iflangja{複製}{replicas}}} {
    \node[cellc, minimum width=\cw mm, fill=ln-box] at (\cx,24) {\ct};
  }
  \foreach \ry/\pa/\pb/\pc in {19/\ttfamily 050012A3/2/{4, 7}, 14/\ttfamily 050012A4/9/---} {
    \node[cellc, minimum width=13mm, fill=white] at (86.5,\ry) {\pa};
    \node[cellc, minimum width=12mm, fill=white] at (99,\ry) {\pb};
    \node[cellc, minimum width=11mm, fill=white] at (110.5,\ry) {\pc};
  }
  \node[lab, text=ln-muted] at (98,8.5) {$\vdots$};
  \draw[->] (70,21) -- (78,21);
\end{tikzpicture}
   (width ~118 mm)
\begin{tikzpicture}[x=1mm, y=1mm, font=\scriptsize\sffamily,
  >={Stealth[length=1.4mm]},
  fld/.style={draw, minimum height=6mm, inner sep=0pt, font=\scriptsize\ttfamily},
  cellc/.style={draw, minimum height=5mm, inner sep=0pt},
  mach/.style={draw, rounded corners=0.8mm, fill=white},
  lab/.style={font=\scriptsize\sffamily, inner sep=0.8pt},
  hd/.style={font=\scriptsize\sffamily\bfseries, inner sep=0.8pt}]
  % page ID
  \node[hd, anchor=south west] at (0,37) {\iflangja{ページ ID}{page ID}};
  \node[fld, fill=ln-accent!20, minimum width=12mm] (nid) at (6,33) {05};
  \node[fld, fill=white, minimum width=22mm] at (23,33) {0012A3};
  \node[lab, anchor=north, text=ln-muted] at (6,29.6) {\iflangja{ノード ID}{node ID}};
  \node[lab, anchor=north, text=ln-muted] at (23,29.6) {PFN};
  % global mapping table
  \node[hd, anchor=south] at (58,37) {\iflangja{グローバルマップ}{global mapping table}};
  \node[lab, anchor=south, text=ln-muted] at (58,34) {\iflangja{（全ノードに複製）}{(replicated on every node)}};
  \node[lab, text=ln-muted] at (51,30.5) {\iflangja{索引}{index}};
  \node[lab, text=ln-muted] at (63,30.5) {\iflangja{管理ノード}{manager}};
  \foreach \ry/\ri/\rm in {26/04/4, 21/05/9, 16/06/6} {
    \node[cellc, minimum width=10mm, fill=white] at (51,\ry) {\ttfamily\ri};
    \node[cellc, minimum width=14mm, fill=white] at (63,\ry) {\rm};
  }
  \node[cellc, minimum width=14mm, fill=ln-accent!20] at (63,21) {9};
  \node[lab, text=ln-muted] at (51,11) {$\vdots$};
  \node[lab, text=ln-muted] at (63,11) {$\vdots$};
  \draw[->] (nid.south) ++(0,-3.6) -- (6,21) -- (46,21);
  % manager node
  \draw[mach] (78,4) rectangle (118,37);
  \node[hd, anchor=north west] at (78.5,36.5) {\iflangja{ノード 9（管理ノード）}{node 9 (manager)}};
  \node[lab, anchor=north west, text=ln-muted] at (78.5,32) {\iflangja{管理するページのメタデータ}{metadata of its pages}};
  \foreach \cx/\cw/\ct in {86.5/13/{\iflangja{ページ}{page}}, 99/12/{\iflangja{所有者}{owner}}, 110.5/11/{\iflangja{複製}{replicas}}} {
    \node[cellc, minimum width=\cw mm, fill=ln-box] at (\cx,24) {\ct};
  }
  \foreach \ry/\pa/\pb/\pc in {19/\ttfamily 050012A3/2/{4, 7}, 14/\ttfamily 050012A4/9/---} {
    \node[cellc, minimum width=13mm, fill=white] at (86.5,\ry) {\pa};
    \node[cellc, minimum width=12mm, fill=white] at (99,\ry) {\pb};
    \node[cellc, minimum width=11mm, fill=white] at (110.5,\ry) {\pc};
  }
  \node[lab, text=ln-muted] at (98,8.5) {$\vdots$};
  \draw[->] (70,21) -- (78,21);
\end{tikzpicture}
   (width ~118 mm)
\begin{tikzpicture}[x=1mm, y=1mm, font=\scriptsize\sffamily,
  >={Stealth[length=1.4mm]},
  fld/.style={draw, minimum height=6mm, inner sep=0pt, font=\scriptsize\ttfamily},
  cellc/.style={draw, minimum height=5mm, inner sep=0pt},
  mach/.style={draw, rounded corners=0.8mm, fill=white},
  lab/.style={font=\scriptsize\sffamily, inner sep=0.8pt},
  hd/.style={font=\scriptsize\sffamily\bfseries, inner sep=0.8pt}]
  % page ID
  \node[hd, anchor=south west] at (0,37) {\iflangja{ページ ID}{page ID}};
  \node[fld, fill=ln-accent!20, minimum width=12mm] (nid) at (6,33) {05};
  \node[fld, fill=white, minimum width=22mm] at (23,33) {0012A3};
  \node[lab, anchor=north, text=ln-muted] at (6,29.6) {\iflangja{ノード ID}{node ID}};
  \node[lab, anchor=north, text=ln-muted] at (23,29.6) {PFN};
  % global mapping table
  \node[hd, anchor=south] at (58,37) {\iflangja{グローバルマップ}{global mapping table}};
  \node[lab, anchor=south, text=ln-muted] at (58,34) {\iflangja{（全ノードに複製）}{(replicated on every node)}};
  \node[lab, text=ln-muted] at (51,30.5) {\iflangja{索引}{index}};
  \node[lab, text=ln-muted] at (63,30.5) {\iflangja{管理ノード}{manager}};
  \foreach \ry/\ri/\rm in {26/04/4, 21/05/9, 16/06/6} {
    \node[cellc, minimum width=10mm, fill=white] at (51,\ry) {\ttfamily\ri};
    \node[cellc, minimum width=14mm, fill=white] at (63,\ry) {\rm};
  }
  \node[cellc, minimum width=14mm, fill=ln-accent!20] at (63,21) {9};
  \node[lab, text=ln-muted] at (51,11) {$\vdots$};
  \node[lab, text=ln-muted] at (63,11) {$\vdots$};
  \draw[->] (nid.south) ++(0,-3.6) -- (6,21) -- (46,21);
  % manager node
  \draw[mach] (78,4) rectangle (118,37);
  \node[hd, anchor=north west] at (78.5,36.5) {\iflangja{ノード 9（管理ノード）}{node 9 (manager)}};
  \node[lab, anchor=north west, text=ln-muted] at (78.5,32) {\iflangja{管理するページのメタデータ}{metadata of its pages}};
  \foreach \cx/\cw/\ct in {86.5/13/{\iflangja{ページ}{page}}, 99/12/{\iflangja{所有者}{owner}}, 110.5/11/{\iflangja{複製}{replicas}}} {
    \node[cellc, minimum width=\cw mm, fill=ln-box] at (\cx,24) {\ct};
  }
  \foreach \ry/\pa/\pb/\pc in {19/\ttfamily 050012A3/2/{4, 7}, 14/\ttfamily 050012A4/9/---} {
    \node[cellc, minimum width=13mm, fill=white] at (86.5,\ry) {\pa};
    \node[cellc, minimum width=12mm, fill=white] at (99,\ry) {\pb};
    \node[cellc, minimum width=11mm, fill=white] at (110.5,\ry) {\pc};
  }
  \node[lab, text=ln-muted] at (98,8.5) {$\vdots$};
  \draw[->] (70,21) -- (78,21);
\end{tikzpicture}

  \caption{分散した状態の索引付け．ページ ID のノード ID フィールドが
    グローバルマッピングテーブルの索引となり，テーブルはページの管理ノードを
    示す．管理ノードは自分が管理するページのメタデータを保持する．エントリ
    \texttt{05} はノード~9 に変更されているので，ノード ID フィールドが
    \texttt{05} であるページは，そのフレームがノード~5 にあるにもかかわらず
    ノード~9 が管理する．}
  \label{fig:l15-indexing}
\end{figure}

\begin{example}
  ある DSM は各ページを 32 ビットのアドレス，すなわち 8 ビットのノード ID
  とそれに続く 24 ビットの PFN で識別する．これにより，最大 256 のノードが
  それぞれ最大 $2^{24}$ 個の \SI{4}{\kibi\byte} フレーム，
  \SI{64}{\gibi\byte} のメモリを提供できる．ページ \texttt{0x050012A3} は
  ノード~5 のフレーム \texttt{0x0012A3} であり，初めはノード~5 がそれを
  管理する．マッピングテーブルはノード ID フィールドで索引され，初めは
  すべての索引 $i$ をノード~$i$ へ写す．ノード~5 が過負荷になると，
  すべてのノードでエントリ~5 を 9 に設定し，ページのメタデータをノード~5 から
  ノード~9 へ移すことによって，ノード~5 のページの管理がノード~9 へ引き渡される
  （\cref{fig:l15-indexing}）．フレーム自体はノード~5 に留まり，ページ ID は
  1つも変わらない．管理ノードをノード ID フィールドから直接取っていたなら，
  これらのページはすべて新しい ID を必要とし，それへのあらゆる参照を更新
  しなければならなかっただろう．
\end{example}

\section{DSM の実装}
\label{sec:l15-implementing}

DSM は2つの理由から共有メモリへのアクセスを捕捉しなければならない．
\source{P4L3，DSM の実装；Nitzberg と Lo；Tanenbaum と Bos ch.~8．}
プロセスが自ノードに
存在しないメモリにアクセスするとき，DSM はそれを保持するノードへメッセージを
送り，アクセスを要求しなければならない．プロセスが，他のノードが複製を保持する
メモリを更新するとき，DSM はその更新を検出し，必要な一貫性操作を行わなければ
ならない．同時に，ローカルに存在し調整を必要としないメモリへのアクセスは，
DSM もオペレーティングシステムも関与させてはならない．あらゆるメモリアクセスに
ソフトウェアが介入すれば，局所性がもたらすはずの性能が失われてしまう．

ページに基づく DSM は，MMU のハードウェア支援によってこの相反する要求を
解決する（\cref{fig:l15-mmu}）．MMU は，仮想アドレスに有効なマッピングが
ないとき，あるいはマッピングが試みられた種類のアクセスを許可しないときは
いつでも，ページフォールト\index{ぺーじふぉーると@ページフォールト}（第8章）
によってオペレーティングシステムへトラップする．それ以外のアクセスはすべて
ハードウェアで完了する．\caution{このページフォールトはスワップ領域から
処理されるのではない．DSM はページを他のノードのメモリから取得する．}DSM は，
捕捉しなければならないアクセスがちょうどフォールトするようにマッピングを
設定する．
\begin{itemize}
  \item 自ノードに存在しないページには有効なマッピングがない．それへの
    アクセスはトラップし，オペレーティングシステムはフォールトを DSM に
    渡す．DSM は所有者にページを要求し，受け取ったコピーのマッピングを設定
    してからプロセスを再開する．読み出しであれば読み出し専用でマップする．
    書き込みであれば，DSM は以下の一貫性操作も行ったうえでページを読み書き可
    でマップし，存在しないページへの書き込みが2回ではなく1回のフォールトで
    済むようにする．
  \item ノードが複製を保持するページは，保護ビット
    \index{ほごびっと@保護ビット}（第8章）によって読み出し専用でマップ
    される．読み出しはトラップせずに進む．書き込みはトラップし，DSM は一貫性
    操作---他の複製を無効化し，所有権を引き受ける---を行ってから書き込みを
    許可し，プロセスを再開する．\sidenote{最初のページに基づく DSM である
    IVY \cite{li1989} はこのように動作し，MRSW アクセスの下で逐次一貫性を
    提供する．後の TreadMarks \cite{keleher1994} は，より弱いモデルの下で，
    変更を加えていない UNIX システム上に，完全にユーザ空間で DSM を実装
    した．}
\end{itemize}
MMU が記録する他の情報も有用である．ページのダーティビット
\index{だーてぃびっと@ダーティビット}は，そのページが変更されたか，したがって
変更を伝播しなければならないかを示す．アクセスビット
\index{あくせすびっと@アクセスビット}は，ローカルな複製が最近使われたかを
示し，メモリが足りなくなったときにどの複製を追い出すかを，ページ置換
（第8章）と同様に決めるのに役立つ．オブジェクトに基づく DSM は，代わりに
言語ランタイムでアクセスを捕捉する．ランタイムがプログラムの共有オブジェクトに
対する操作を仲介するからである．\sidenote{これらの MMU の操作はただでは
ない．Appel と Li \cite{appel1991} の測定では，1991 年当時の平均的なシステムは
フォールトと保護の変更に \SI{1200}{\micro\second} ほどを費やし，さらに
マッピングを保持する他の CPU で TLB シュートダウンを要した．}

\begin{figure}[htbp]
  \centering
  % Implementing a page-based DSM with the MMU: accesses permitted by the
% current mapping complete in hardware; the others fault, and the DSM fetches
% the page, performs consistency operations, or both, before the access is
% repeated.
% Usage: % Implementing a page-based DSM with the MMU: accesses permitted by the
% current mapping complete in hardware; the others fault, and the DSM fetches
% the page, performs consistency operations, or both, before the access is
% repeated.
% Usage: % Implementing a page-based DSM with the MMU: accesses permitted by the
% current mapping complete in hardware; the others fault, and the DSM fetches
% the page, performs consistency operations, or both, before the access is
% repeated.
% Usage: \input{figures/l15-mmu}   (width ~112 mm)
\begin{tikzpicture}[x=1mm, y=1mm, font=\scriptsize\sffamily,
  >={Stealth[length=1.5mm]},
  box/.style={draw, fill=white, align=center, inner sep=1.2pt, minimum height=8mm, text width=30mm},
  dec/.style={draw, diamond, aspect=2.2, fill=ln-box, align=center, inner sep=0pt},
  lab/.style={font=\scriptsize\sffamily, inner sep=0.8pt}]
  \node[box] (acc) at (17,50) {\iflangja{プロセスが共有アドレス\\にアクセス}{process accesses a\\shared address}};
  \node[dec] (mmu) at (17,33) {\iflangja{MMU：アクセス\\は許可？}{MMU: access\\permitted?}};
  \node[box, fill=ln-tip!12] (ok) at (17,14) {\iflangja{ハードウェアで完了\\（トラップなし）}{access completes in\\hardware (no trap)}};
  \node[box, fill=ln-source!12, text width=34mm] (trap) at (66,33) {\iflangja{ページフォールト：\\OS へトラップ，\\OS は DSM に渡す}{page fault: trap to the OS,\\which hands it to the DSM}};
  \node[box] (fetch) at (50,10) {\iflangja{所有者にページを要求：\\読みなら読み出し専用\\でマップ}{request the page from its\\owner; map it read-only\\for a read}};
  \node[box, text width=32mm] (coord) at (93,10) {\iflangja{一貫性操作：他の複製を\\無効化，読み書き可でマップ}{consistency operations:\\invalidate other replicas;\\map it read-write}};
  \draw[->] (acc) -- (mmu);
  \draw[->] (mmu) -- node[lab, right] {\iflangja{はい}{yes}} (ok);
  \draw[->] (mmu) -- node[lab, above] {\iflangja{いいえ}{no}} (trap);
  \draw[->] (trap.south) -- ++(0,-5) -| (fetch.north);
  \draw[->] (trap.south) -- ++(0,-5) -| (coord.north);
  \node[lab, anchor=east, text=ln-muted, align=right] at (49.2,20) {\iflangja{ローカルに\\ない}{no local\\copy}};
  \draw[->] (fetch.east) -- node[lab, above, inner sep=0.6pt] {\iflangja{書き}{write}} (coord.west);
  \node[lab, anchor=west, text=ln-muted] at (50.6,4.2) {\iflangja{読み}{read}};
  \node[lab, anchor=west, text=ln-muted, align=left] at (93.8,20) {\iflangja{読み出し専用\\への書き込み}{write to a\\read-only copy}};
  \draw[->] (fetch.south) -- ++(0,-5) -| (ok.south);
  \draw (coord.south) |- ++(-40,-5);
  \node[lab, anchor=north, text=ln-muted] at (73,-0.4) {\iflangja{プロセスを再開：アクセスを再実行}{resume the process: the access is repeated}};
\end{tikzpicture}
   (width ~112 mm)
\begin{tikzpicture}[x=1mm, y=1mm, font=\scriptsize\sffamily,
  >={Stealth[length=1.5mm]},
  box/.style={draw, fill=white, align=center, inner sep=1.2pt, minimum height=8mm, text width=30mm},
  dec/.style={draw, diamond, aspect=2.2, fill=ln-box, align=center, inner sep=0pt},
  lab/.style={font=\scriptsize\sffamily, inner sep=0.8pt}]
  \node[box] (acc) at (17,50) {\iflangja{プロセスが共有アドレス\\にアクセス}{process accesses a\\shared address}};
  \node[dec] (mmu) at (17,33) {\iflangja{MMU：アクセス\\は許可？}{MMU: access\\permitted?}};
  \node[box, fill=ln-tip!12] (ok) at (17,14) {\iflangja{ハードウェアで完了\\（トラップなし）}{access completes in\\hardware (no trap)}};
  \node[box, fill=ln-source!12, text width=34mm] (trap) at (66,33) {\iflangja{ページフォールト：\\OS へトラップ，\\OS は DSM に渡す}{page fault: trap to the OS,\\which hands it to the DSM}};
  \node[box] (fetch) at (50,10) {\iflangja{所有者にページを要求：\\読みなら読み出し専用\\でマップ}{request the page from its\\owner; map it read-only\\for a read}};
  \node[box, text width=32mm] (coord) at (93,10) {\iflangja{一貫性操作：他の複製を\\無効化，読み書き可でマップ}{consistency operations:\\invalidate other replicas;\\map it read-write}};
  \draw[->] (acc) -- (mmu);
  \draw[->] (mmu) -- node[lab, right] {\iflangja{はい}{yes}} (ok);
  \draw[->] (mmu) -- node[lab, above] {\iflangja{いいえ}{no}} (trap);
  \draw[->] (trap.south) -- ++(0,-5) -| (fetch.north);
  \draw[->] (trap.south) -- ++(0,-5) -| (coord.north);
  \node[lab, anchor=east, text=ln-muted, align=right] at (49.2,20) {\iflangja{ローカルに\\ない}{no local\\copy}};
  \draw[->] (fetch.east) -- node[lab, above, inner sep=0.6pt] {\iflangja{書き}{write}} (coord.west);
  \node[lab, anchor=west, text=ln-muted] at (50.6,4.2) {\iflangja{読み}{read}};
  \node[lab, anchor=west, text=ln-muted, align=left] at (93.8,20) {\iflangja{読み出し専用\\への書き込み}{write to a\\read-only copy}};
  \draw[->] (fetch.south) -- ++(0,-5) -| (ok.south);
  \draw (coord.south) |- ++(-40,-5);
  \node[lab, anchor=north, text=ln-muted] at (73,-0.4) {\iflangja{プロセスを再開：アクセスを再実行}{resume the process: the access is repeated}};
\end{tikzpicture}
   (width ~112 mm)
\begin{tikzpicture}[x=1mm, y=1mm, font=\scriptsize\sffamily,
  >={Stealth[length=1.5mm]},
  box/.style={draw, fill=white, align=center, inner sep=1.2pt, minimum height=8mm, text width=30mm},
  dec/.style={draw, diamond, aspect=2.2, fill=ln-box, align=center, inner sep=0pt},
  lab/.style={font=\scriptsize\sffamily, inner sep=0.8pt}]
  \node[box] (acc) at (17,50) {\iflangja{プロセスが共有アドレス\\にアクセス}{process accesses a\\shared address}};
  \node[dec] (mmu) at (17,33) {\iflangja{MMU：アクセス\\は許可？}{MMU: access\\permitted?}};
  \node[box, fill=ln-tip!12] (ok) at (17,14) {\iflangja{ハードウェアで完了\\（トラップなし）}{access completes in\\hardware (no trap)}};
  \node[box, fill=ln-source!12, text width=34mm] (trap) at (66,33) {\iflangja{ページフォールト：\\OS へトラップ，\\OS は DSM に渡す}{page fault: trap to the OS,\\which hands it to the DSM}};
  \node[box] (fetch) at (50,10) {\iflangja{所有者にページを要求：\\読みなら読み出し専用\\でマップ}{request the page from its\\owner; map it read-only\\for a read}};
  \node[box, text width=32mm] (coord) at (93,10) {\iflangja{一貫性操作：他の複製を\\無効化，読み書き可でマップ}{consistency operations:\\invalidate other replicas;\\map it read-write}};
  \draw[->] (acc) -- (mmu);
  \draw[->] (mmu) -- node[lab, right] {\iflangja{はい}{yes}} (ok);
  \draw[->] (mmu) -- node[lab, above] {\iflangja{いいえ}{no}} (trap);
  \draw[->] (trap.south) -- ++(0,-5) -| (fetch.north);
  \draw[->] (trap.south) -- ++(0,-5) -| (coord.north);
  \node[lab, anchor=east, text=ln-muted, align=right] at (49.2,20) {\iflangja{ローカルに\\ない}{no local\\copy}};
  \draw[->] (fetch.east) -- node[lab, above, inner sep=0.6pt] {\iflangja{書き}{write}} (coord.west);
  \node[lab, anchor=west, text=ln-muted] at (50.6,4.2) {\iflangja{読み}{read}};
  \node[lab, anchor=west, text=ln-muted, align=left] at (93.8,20) {\iflangja{読み出し専用\\への書き込み}{write to a\\read-only copy}};
  \draw[->] (fetch.south) -- ++(0,-5) -| (ok.south);
  \draw (coord.south) |- ++(-40,-5);
  \node[lab, anchor=north, text=ln-muted] at (73,-0.4) {\iflangja{プロセスを再開：アクセスを再実行}{resume the process: the access is repeated}};
\end{tikzpicture}

  \caption{ページに基づく DSM におけるアクセスの捕捉．現在のマッピングが
    許可するアクセスはハードウェアから出ない．それ以外はフォールトし，DSM は
    ページを取得するか一貫性操作を行い，その後アクセスがやり直される．}
  \label{fig:l15-mmu}
\end{figure}

\begin{example}
  ページ $p$ は初め，所有者であるノード~1 にのみ存在し，ノード~1 はそれを
  読み書き可でマップしている．\Cref{tab:l15-mmu} は，書き込み時に無効化する
  MRSW アクセスの下での3つのノードによるアクセスを追ったものである．
  ステップ~1 でノード~2 が複製を得ると，ノード~1 のマッピングは読み出し専用に
  格下げされ，後にノード~1 が書き込めばフォールトして複製を無効化するように
  なる．7 回のアクセスのうち 4 回がフォールトする．ステップ~2 と~3 の読み出し，
  およびステップ~5 の書き込みは，マッピングがすでにそれらを許可しているので
  ハードウェアで完了する．\tip{ノードごとに見れば，ページのマッピングは
  「なし，RO，RW」の3状態マシンであり，フォールトするのは遷移だけである．
  読み出しはなしから RO へ，書き込みは RO から RW へ移し，他のノードの
  アクセスは RO またはなしへ格下げする．}
\end{example}

\begin{table}[htbp]
  \centering
  \caption{ページ $p$ へのアクセスと，それが引き起こすフォールト．最後の列は
    各ステップの後のノード 1/2/3 における $p$ のマッピングを示す（RW は読み書き
    可，RO は読み出し専用，--- はマッピングなし）．}
  \label{tab:l15-mmu}
  \small
  \setlength{\tabcolsep}{4pt}
  \begin{tabular}{@{}rll>{\raggedright\arraybackslash}p{0.29\linewidth}l@{}}
    \toprule
    ステップ & アクセス & フォールト & DSM の動作 & マッピング \\
    \midrule
    0 & （初期状態） & --- & --- & RW/---/--- \\
    1 & ノード 2 が読む & ページフォールト & 所有者のノード~1 から $p$ を
      取得；ノード~1 を RO に格下げ & RO/RO/--- \\
    2 & ノード 2 が読む & なし & --- & RO/RO/--- \\
    3 & ノード 1 が読む & なし & --- & RO/RO/--- \\
    4 & ノード 2 が書く & 保護違反 & ノード~1 のコピーを無効化；ノード~2 が
      所有者になる & ---/RW/--- \\
    5 & ノード 2 が書く & なし & --- & ---/RW/--- \\
    6 & ノード 3 が読む & ページフォールト & 所有者のノード~2 から $p$ を
      取得；ノード~2 を RO に格下げ & ---/RO/RO \\
    7 & ノード 2 が書く & 保護違反 & ノード~3 のコピーを無効化 &
      ---/RW/--- \\
    \bottomrule
  \end{tabular}
\end{table}

\begin{keypoint}
  ページに基づく DSM は，共有されるページをすべてのノードが提供したメモリに
  置き，複製されたグローバルマップを通じて各ページの管理ノードを見つけ，MMU を
  使って，必要なアクセスだけ---存在しないページへのアクセスと，複製された
  ページへの書き込み---をちょうど捕捉する．それ以外のアクセスはすべて
  ハードウェアの速度で進む．
\end{keypoint}

\needspace{6\baselineskip}
\section{一貫性モデル}
\label{sec:l15-models}

\cref{sec:l15-management} の機構は異なる時点で起動でき，その結果得られる
保証も異なる．\source{P4L3，一貫性モデル，厳密一貫性，逐次一貫性，因果
一貫性，弱一貫性；Tanenbaum と van Steen \cite{tanenbaum2007} ch.~7．}
分散ファイルシステムでは，これらの保証は第14章のファイル共有セマンティクスで
ある．メモリについて，そしてより一般に共有状態について，それを述べるのが
一貫性モデルである．

\begin{definition}[一貫性モデル]
  \term[いっかんせいもでる]{一貫性モデル}[consistency model]とは，メモリ---
  あるいは任意の共有状態---と，その上のソフトウェア層との間の取り決めで
  ある．ソフトウェアが特定の規則に従う限り，メモリは，アクセスの順序づけと，
  更新の伝播および可視性に関して正しく振る舞うことを保証する．
\end{definition}

したがってファイル共有セマンティクスは，ファイルに対する一貫性モデルである．
メモリに対する一貫性モデルを記述するために，メモリ位置 $m$ から値 $v$ を
返す読み出しを $R_{m}(v)$，$m$ への $v$ の書き込みを $W_{m}(v)$ と書く．
操作はプロセスごとに1本のタイムラインに描き，時間は右へ進む．すべての位置は
初め 0 を保持している．\margin{タイムラインにおいて $m_a$ と $m_b$ は共有
アドレス空間の2つの位置であり，$t_1 < t_2 < \dots$ は外部の観測者が見る
実時刻の時点である．}

\subsection{厳密一貫性}

\begin{definition}[厳密一貫性]
  \term[げんみついっかんせい]{厳密一貫性}[strict consistency]の下では，
  あらゆる更新がただちにどこからも見える．すなわち，どのノードのどの
  読み出しも，その位置に対する実時間で最も新しい書き込みの値を返し，
  すべてのノードが更新を，それがなされた順序で観測する．
\end{definition}

厳密一貫性は，プログラマがメモリに対して直観的に期待するものだが，理論上の
モデルである．単一の共有メモリ型マルチプロセッサ上でさえ，異なる CPU 上の
並行な更新は，ソフトウェアがロックと同期を加えない限り，保証された順序で
効果を持つわけではない．分散システムでは，更新はゼロでないレイテンシを持つ
ネットワークを渡り，メッセージは順序が入れ替わったり失われたりし，異なる
ノード上の2つの書き込みのどちらがより新しいかを定めるグローバルな時計も
存在しない．厳密一貫性を保証することはできない．

\subsection{逐次一貫性}

\begin{definition}[逐次一貫性]
  \term[ちくじいっかんせい]{逐次一貫性}[sequential consistency]の下では，
  あらゆる実行の結果が，すべてのプロセスの操作をある逐次順序で実行したときと
  同じになる．その順序では，各プロセスの操作はそのプロセスが発行した順に
  現れ，どの読み出しもその位置に対する最も新しい書き込みを返し，すべての
  プロセスがこの同一の順序を観測する \cite{lamport1979}．
\end{definition}

したがって，異なるプロセスからの更新は任意にインタリーブされてよく，
インタリーブが実時間の順序と一致する必要もないが，すべてのプロセスが同じ
インタリーブを見なければならない．\caution{逐次一貫な実行は実時間と一致する
必要はない．すべてのプロセスが同じ順序を観測する限り，読み出しが，実時間で
それに最も近く先行する書き込みより古い値を返してもよい．}

\subsection{因果一貫性}

ある書き込みが以前の書き込みに\emph{因果的に依存する}とは，それを発行した
プロセスが，直接に---たとえばその書き込んだ値を読むことによって---あるいは
そうした依存の連鎖を通じて，以前の書き込みを観測していたことをいう．どちらも
他方に依存しない書き込みは\emph{並行}であるという．

\begin{definition}[因果一貫性]
  \term[いんがいっかんせい]{因果一貫性}[causal consistency]の下では，
  すべてのプロセスが，因果的に関係する書き込みをその因果順序で観測し，各
  プロセスの操作をそのプロセスが発行した順序で観測する．並行な書き込みは，
  プロセスごとに異なる順序で観測されてよい．
\end{definition}

因果一貫性を提供するには，DSM は因果関係がありうる箇所を検出しなければ
ならない．すなわち，各書き込みの前にそのプロセスがどの更新を観測していたかを
記録しなければならない．\margin{ベクタタイムスタンプはこうした依存を記録する
一般的な方法である（Tanenbaum と van Steen ch.~6）．}その代わり，順序を
調整する必要があるのは因果的に関係する更新だけである．

\subsection{弱一貫性}

弱一貫性は\emph{同期点}を導入する．同期点とは，DSM がその上のソフトウェアに
提供する \texttt{sync} のような操作である．

\begin{definition}[弱一貫性]
  \term[じゃくいっかんせい]{弱一貫性}[weak consistency]の下では，同期点より
  前になされたあらゆる更新が，同期点より後にはすべてのノードから見える．
  同期点の間では，ノードがどの更新を観測するかについて何の保証も与えられ
  ない．
\end{definition}

DSM は，共有状態全体に対して1つの同期操作を提供することも，個々のページや
オブジェクトといった状態の部分集合に対して別々の操作を提供し，アプリケー
ションが必要な状態だけを同期できるようにすることもできる．また操作を2つに
分けることもできる．プロセスが共有状態を使う前に行い，他のプロセスの更新を
そのプロセスから見えるようにする \emph{acquire} と，更新を終えた後に行い，
自身の更新を伝播する \emph{release} である．\margin{セッションセマンティクス
（第14章）は弱一貫性のファイル版である．\texttt{open} と \texttt{close} が
同期点にあたる．}弱一貫性は，データの移動と調整を同期点に限り，そこでは前の
同期点以降になされた更新をまとめて伝播できる．その代わり，DSM はこれらの更新を
記録し集約するための追加の状態を保持しなければならず，同期点の間ではプロセスが
古いデータを観測しうる．

\subsection{モデルの比較}

厳密一貫性，逐次一貫性，因果一貫性は階層をなす．厳密一貫な実行はすべて逐次
一貫であり，逐次一貫な実行はすべて因果一貫である．より弱いモデルはより多くの
実行を許すので調整が少なくて済むが，アプリケーションに与える保証は少ない．
\caution{実際のハードウェアも逐次一貫ではない．x86 では，ロードが別の位置への
先行するストアを追い越しうる（total store order）．ARM や POWER はさらに順序を
入れ替える．フェンスが同期点の役割を果たす（Intel SDM vol.~3A，メモリ順序
付け）．}

\begin{figure}[htbp]
  \centering
  % Operations for the consistency-model example: P1 and P2 write m_a and m_b,
% P3 and P4 read both locations in opposite orders, a synchronization point
% follows at t5, and P4 reads m_a again at t6.  P2's read of m_a (grey) is
% present only in the variant with a causal dependency.
% Usage: % Operations for the consistency-model example: P1 and P2 write m_a and m_b,
% P3 and P4 read both locations in opposite orders, a synchronization point
% follows at t5, and P4 reads m_a again at t6.  P2's read of m_a (grey) is
% present only in the variant with a causal dependency.
% Usage: % Operations for the consistency-model example: P1 and P2 write m_a and m_b,
% P3 and P4 read both locations in opposite orders, a synchronization point
% follows at t5, and P4 reads m_a again at t6.  P2's read of m_a (grey) is
% present only in the variant with a causal dependency.
% Usage: \input{figures/l15-models}   (width ~112 mm)
\begin{tikzpicture}[x=1mm, y=1mm, font=\scriptsize\sffamily,
  row/.style={font=\scriptsize\sffamily, anchor=east},
  lab/.style={font=\scriptsize, inner sep=0.8pt},
  ev/.style={circle, fill=black, inner sep=0pt, minimum size=1.3mm}]
  % t_i -> x = 8 + 16 i (mm); rows: P1 27, P2 18, P3 9, P4 0
  \foreach \pn/\py in {1/27,2/18,3/9,4/0} {
    \node[row] at (14,\py) {P\pn};
    \draw[ln-rule] (16,\py) -- (112,\py);
  }
  % synchronization point
  \draw[densely dashed, ln-accent, thick] (88,-3) -- (88,31);
  \node[lab, anchor=south, text=ln-accent, font=\scriptsize\sffamily] at (88,31) {sync};
  % P1, P2 writes
  \node[ev] at (24,27) {};  \node[lab, anchor=south] at (24,28) {$W_{m_a}(7)$};
  \node[ev] at (40,18) {};  \node[lab, anchor=south] at (40,19) {$W_{m_b}(9)$};
  \node[ev, fill=ln-muted] at (31,18) {};
  \node[lab, anchor=north, text=ln-muted] at (29,17) {($R_{m_a}(7)$)};
  % P3 reads
  \node[ev] at (56,9) {};  \node[lab, anchor=south] at (56,10) {$R_{m_a}(?)$};
  \node[ev] at (72,9) {};  \node[lab, anchor=south] at (72,10) {$R_{m_b}(?)$};
  % P4 reads
  \node[ev] at (56,0) {};  \node[lab, anchor=south] at (56,1) {$R_{m_b}(?)$};
  \node[ev] at (72,0) {};  \node[lab, anchor=south] at (72,1) {$R_{m_a}(?)$};
  \node[ev] at (104,0) {}; \node[lab, anchor=south] at (104,1) {$R_{m_a}(?)$};
  % time axis
  \draw[->] (16,-6) -- (113,-6);
  \foreach \ti in {1,...,6} {
    \draw ({8+16*\ti},-5.2) -- ({8+16*\ti},-6.8);
    \node[lab, anchor=north] at ({8+16*\ti},-7) {$t_{\ti}$};
  }
\end{tikzpicture}
   (width ~112 mm)
\begin{tikzpicture}[x=1mm, y=1mm, font=\scriptsize\sffamily,
  row/.style={font=\scriptsize\sffamily, anchor=east},
  lab/.style={font=\scriptsize, inner sep=0.8pt},
  ev/.style={circle, fill=black, inner sep=0pt, minimum size=1.3mm}]
  % t_i -> x = 8 + 16 i (mm); rows: P1 27, P2 18, P3 9, P4 0
  \foreach \pn/\py in {1/27,2/18,3/9,4/0} {
    \node[row] at (14,\py) {P\pn};
    \draw[ln-rule] (16,\py) -- (112,\py);
  }
  % synchronization point
  \draw[densely dashed, ln-accent, thick] (88,-3) -- (88,31);
  \node[lab, anchor=south, text=ln-accent, font=\scriptsize\sffamily] at (88,31) {sync};
  % P1, P2 writes
  \node[ev] at (24,27) {};  \node[lab, anchor=south] at (24,28) {$W_{m_a}(7)$};
  \node[ev] at (40,18) {};  \node[lab, anchor=south] at (40,19) {$W_{m_b}(9)$};
  \node[ev, fill=ln-muted] at (31,18) {};
  \node[lab, anchor=north, text=ln-muted] at (29,17) {($R_{m_a}(7)$)};
  % P3 reads
  \node[ev] at (56,9) {};  \node[lab, anchor=south] at (56,10) {$R_{m_a}(?)$};
  \node[ev] at (72,9) {};  \node[lab, anchor=south] at (72,10) {$R_{m_b}(?)$};
  % P4 reads
  \node[ev] at (56,0) {};  \node[lab, anchor=south] at (56,1) {$R_{m_b}(?)$};
  \node[ev] at (72,0) {};  \node[lab, anchor=south] at (72,1) {$R_{m_a}(?)$};
  \node[ev] at (104,0) {}; \node[lab, anchor=south] at (104,1) {$R_{m_a}(?)$};
  % time axis
  \draw[->] (16,-6) -- (113,-6);
  \foreach \ti in {1,...,6} {
    \draw ({8+16*\ti},-5.2) -- ({8+16*\ti},-6.8);
    \node[lab, anchor=north] at ({8+16*\ti},-7) {$t_{\ti}$};
  }
\end{tikzpicture}
   (width ~112 mm)
\begin{tikzpicture}[x=1mm, y=1mm, font=\scriptsize\sffamily,
  row/.style={font=\scriptsize\sffamily, anchor=east},
  lab/.style={font=\scriptsize, inner sep=0.8pt},
  ev/.style={circle, fill=black, inner sep=0pt, minimum size=1.3mm}]
  % t_i -> x = 8 + 16 i (mm); rows: P1 27, P2 18, P3 9, P4 0
  \foreach \pn/\py in {1/27,2/18,3/9,4/0} {
    \node[row] at (14,\py) {P\pn};
    \draw[ln-rule] (16,\py) -- (112,\py);
  }
  % synchronization point
  \draw[densely dashed, ln-accent, thick] (88,-3) -- (88,31);
  \node[lab, anchor=south, text=ln-accent, font=\scriptsize\sffamily] at (88,31) {sync};
  % P1, P2 writes
  \node[ev] at (24,27) {};  \node[lab, anchor=south] at (24,28) {$W_{m_a}(7)$};
  \node[ev] at (40,18) {};  \node[lab, anchor=south] at (40,19) {$W_{m_b}(9)$};
  \node[ev, fill=ln-muted] at (31,18) {};
  \node[lab, anchor=north, text=ln-muted] at (29,17) {($R_{m_a}(7)$)};
  % P3 reads
  \node[ev] at (56,9) {};  \node[lab, anchor=south] at (56,10) {$R_{m_a}(?)$};
  \node[ev] at (72,9) {};  \node[lab, anchor=south] at (72,10) {$R_{m_b}(?)$};
  % P4 reads
  \node[ev] at (56,0) {};  \node[lab, anchor=south] at (56,1) {$R_{m_b}(?)$};
  \node[ev] at (72,0) {};  \node[lab, anchor=south] at (72,1) {$R_{m_a}(?)$};
  \node[ev] at (104,0) {}; \node[lab, anchor=south] at (104,1) {$R_{m_a}(?)$};
  % time axis
  \draw[->] (16,-6) -- (113,-6);
  \foreach \ti in {1,...,6} {
    \draw ({8+16*\ti},-5.2) -- ({8+16*\ti},-6.8);
    \node[lab, anchor=north] at ({8+16*\ti},-7) {$t_{\ti}$};
  }
\end{tikzpicture}

  \caption{\cref{ex:l15-models} の操作．P2 による $m_a$ の読み出し（灰色）は，
    P2 の書き込みが P1 の書き込みに依存する変種でのみ起こる．破線は同期点で
    ある．}
  \label{fig:l15-models}
\end{figure}

\begin{example}
  \label{ex:l15-models}
  \cref{fig:l15-models} では，P1 が $t_1$ に $m_a$ へ 7 を書き，P2 が $t_2$
  に $m_b$ へ 9 を書く．P3 と P4 は $t_3$ と $t_4$ に両方の位置を，互いに
  逆の順序で読む．\Cref{tab:l15-models} は，これらの読み出しの3通りの結果と，
  それが満たすモデルを挙げたものである．
  \begin{itemize}
    \item 結果 (a) は厳密一貫である．どの読み出しも実時間で最も新しい
      書き込みを返している．
    \item 結果 (b) はそうではない．P4 は，$t_2$ に 9 が書かれた後の $t_3$ に
      $m_b$ から 0 を読んでいる．しかし逐次一貫ではある．P3 の
      $R_{m_a}(0)$，P4 の $R_{m_b}(0)$，$W_{m_a}(7)$，$W_{m_b}(9)$，P3 の
      $R_{m_b}(9)$，P4 の $R_{m_a}(7)$ という順序が，どの読み出しも説明し，
      各プロセスの操作の順序も保っているからである．\tip{逐次一貫性を
      確かめるには，各プロセスの順序を保ったまま全操作を1本の列にまとめ，
      すべての読み出しを説明できるものを探せばよい．}
    \item 結果 (c) は逐次一貫ではない．P3 は $W_{m_a}(7)$ を $W_{m_b}(9)$
      より先に観測し，P4 は逆の順序で観測しており，両者を説明する単一の順序は
      存在しない．2つの書き込みが並行であれば，(c) は因果一貫である．しかし
      P2 が $m_b$ に書く前に $m_a$ から 7 を読んでいたなら，$W_{m_b}(9)$ は
      $W_{m_a}(7)$ に因果的に依存し，$m_b$ から 9 を読みその後 $m_a$ から
      0 を読む P4 は，原因より先に結果を観測していることになる．
  \end{itemize}
  $t_5$ に同期点があれば，3つの結果はいずれも弱一貫である．$t_5$ より前には
  何の保証も適用されないからである．ただし $t_6$ における P4 の $m_a$ の
  読み出しは 7 を返さなければならない．\caution{モデルは結果を制約するので
  あって，結果を決定するのではない．逐次一貫性を提供する DSM は，ある実行で
  (a) を，次の実行で (b) を生じうるし，弱一貫な DSM は3つのいずれをも
  生じうる．}
\end{example}

\begin{table}[htbp]
  \centering
  \caption{\cref{ex:l15-models} における $t_3$ と $t_4$ の読み出しの結果．
    因果の列は，書き込みが並行である場合／P2 が $m_b$ に書く前に $m_a$ を
    読んだ場合の結果を示す．}
  \label{tab:l15-models}
  \small
  \begin{tabular}{@{}lllcccc@{}}
    \toprule
    & P3 & P4 & 厳密 & 逐次 & 因果 & 弱 \\
    \midrule
    (a) & $R_{m_a}(7)$, $R_{m_b}(9)$ & $R_{m_b}(9)$, $R_{m_a}(7)$ & はい
      & はい & はい／はい & はい \\
    (b) & $R_{m_a}(0)$, $R_{m_b}(9)$ & $R_{m_b}(0)$, $R_{m_a}(7)$ & いいえ
      & はい & はい／はい & はい \\
    (c) & $R_{m_a}(7)$, $R_{m_b}(0)$ & $R_{m_b}(9)$, $R_{m_a}(0)$ & いいえ
      & いいえ & はい／いいえ & はい \\
    \bottomrule
  \end{tabular}
\end{table}

\begin{keypoint}[章のまとめ]
  分散共有メモリは，どのノードもメモリを提供し，それへのアクセスを処理し，
  一貫性プロトコルに参加するピア分散アプリケーションであり，それによって
  多数のマシン上のプロセスが1つのアドレス空間を共有する．その設計は，
  共有の粒度（ページかオブジェクトか，偽共有の危険を伴う），アクセス
  アルゴリズム（SRSW，MRSW，MRMW），局所性を得るためのマイグレーションと
  レプリケーションの選択，そして無効化をプッシュするか変更をプルするかと
  いう一貫性管理によって決まる．ページに基づく DSM は，複製されたマッピング
  テーブルを通じて各ページの管理ノードを見つけ，調整を必要とするアクセス
  だけを捕捉するために MMU のフォールトに依存する．一貫性モデルはプロセスが
  何を観測するかを述べる．厳密一貫性は実際には達成できない．逐次一貫性は
  すべての操作に共通の1つの順序を要求する．因果一貫性は因果的に関係する
  書き込みだけを順序づける．弱一貫性は同期点においてのみ可視性を保証する．
\end{keypoint}

\end{document}
